% !TeX program = xelatex
\documentclass[11pt, a4paper]{report}
%----------------------------------------------------------------------------------------
%	PACKAGES (宏包管理)
%----------------------------------------------------------------------------------------
% 1. 编码与语言
\usepackage[english]{babel}
\usepackage{fontspec}
\usepackage{xeCJK}
\setmainfont{Latin Modern Roman}[
    SmallCapsFont={Latin Modern Roman Caps}
]
\setCJKmainfont{Microsoft YaHei}[AutoFakeSlant=0.2, AutoFakeBold=2]
% 2. 页面布局与排版
\usepackage{geometry}
\geometry{
    textwidth=6in, % 稍微加宽以便放入更好的box
    top=1in,
    bottom=1in,
    headheight=12pt,
    headsep=25pt,
    footskip=30pt
}
\usepackage{setspace}
\setstretch{1.2}
\usepackage{float}
\usepackage[shortlabels]{enumitem} % 更好的列表控制
\usepackage{tocloft}  % 加载目录控制宏包
\renewcommand{\cftsecleader}{\cftdotfill{\cftdotsep}} % 强制让 section 后面显示点号
% 3. 数学支持
\usepackage{amsmath, amsthm, amssymb, amsfonts, amscd}
\usepackage{mathtools}
\usepackage{mathrsfs}
\usepackage{thmtools} % 定理环境增强
\usepackage{extarrows}
% 4. 图形与颜色
\usepackage{graphicx}
\usepackage[dvipsnames, svgnames]{xcolor} % 加载更多颜色名
\usepackage{tikz}
\usepackage{tikz-cd}
\usepackage{ytableau} % 杨表
\usetikzlibrary{decorations.pathreplacing, calc, shadows}
% 5. 漂亮的文本框
\usepackage{tcolorbox}
\tcbuselibrary{breakable, skins, theorems}
% 6. 参考文献
%\usepackage[backend=biber,style=gb7714-2015,sorting=gb7714-2015]{biblatex}
%\addbibresource{references.bib}
% 7. 交互 (最后加载)
\usepackage[hidelinks]{hyperref}
\usepackage{array}
%----------------------------------------------------------------------------------------
%	CUSTOM COLORS (自定义配色方案)
%----------------------------------------------------------------------------------------
% 定理色系 - 蓝色
\definecolor{ThmColor}{RGB}{0, 71, 171}        % 深钴蓝
\definecolor{ThmBg}{RGB}{240, 244, 255}        % 极浅蓝
% 定义色系 - 绿色
\definecolor{DefColor}{RGB}{0, 100, 0}         % 深绿
\definecolor{DefBg}{RGB}{240, 255, 240}        % 蜜瓜绿
% 引理/命题色系 - 橙色
\definecolor{LemColor}{RGB}{204, 85, 0}        % 焦橙色
\definecolor{LemBg}{RGB}{255, 248, 240}        % 极浅橙
% 示例/注记色系 - 灰色/紫色
\definecolor{ExColor}{RGB}{105, 105, 105}      % 暗灰
\definecolor{ExBg}{RGB}{250, 250, 250}         % 极浅灰
% 专题/话题色系 - 靛紫色
\definecolor{TopicColor}{RGB}{75, 0, 130}      % 靛蓝色 (Indigo)
\definecolor{TopicBg}{RGB}{248, 245, 255}      % 极浅紫
%----------------------------------------------------------------------------------------
%	THEOREM ENVIRONMENTS (环境美化核心)
%----------------------------------------------------------------------------------------
% 1. Theorem (蓝色风格)
\declaretheorem[name=Theorem, numberwithin=section]{theorem}
\tcolorboxenvironment{theorem}{
    enhanced jigsaw,
    colframe=ThmColor,       % 边框颜色
    colback=ThmBg,           % 背景颜色
    coltitle=white,          
    boxrule=0.5pt,           % 边框粗细
    leftrule=2pt,            % 左边框加粗
    sharp corners=downhill,  % 样式
    breakable,               % 允许跨页
    drop shadow=black!5!white % 轻微阴影
}
% Corollary 使用与 Theorem 相同的风格，但通常紧随其后
\declaretheorem[name=Corollary, numberlike=theorem]{corollary}
\tcolorboxenvironment{corollary}{
    enhanced jigsaw,
    colframe=ThmColor!80!white, 
    colback=ThmBg,
    boxrule=0pt, leftrule=2pt, arc=0pt, outer arc=0pt,
    breakable
}
% 2. Definition (绿色风格)
\declaretheorem[name=Definition, numberlike=theorem]{definition}
\tcolorboxenvironment{definition}{
    enhanced jigsaw,
    colframe=DefColor,
    colback=DefBg,
    boxrule=0.5pt, leftrule=2pt,
    sharp corners=downhill,
    breakable,
    drop shadow=black!5!white
}
% 3. Lemma & Proposition (橙色风格)
\declaretheorem[name=Lemma, numberlike=theorem]{lemma}
\declaretheorem[name=Proposition, numberlike=theorem]{proposition}
\tcolorboxenvironment{lemma}{
    enhanced jigsaw,
    colframe=LemColor,
    colback=LemBg,
    boxrule=0.5pt, leftrule=2pt,
    breakable
}
\tcolorboxenvironment{proposition}{
    enhanced jigsaw,
    colframe=LemColor,
    colback=LemBg,
    boxrule=0.5pt, leftrule=2pt,
    breakable
}
% 4. Example & Remark (简约风格)
\declaretheorem[name=Example, numberlike=theorem]{example}
\declaretheorem[name=Remark, numberlike=theorem]{remark}
\declaretheorem[name=Exercise, numberlike=theorem]{exercise}
\tcolorboxenvironment{example}{
    enhanced jigsaw,
    colframe=ExColor,
    colback=ExBg,
    leftrule=2pt, boxrule=0pt,
    frame hidden, % 隐藏其他三边
    breakable
}
% 5. Topic (靛紫色风格，带标题支持)
\declaretheorem[name=Topic, numberlike=theorem]{topic}
\tcolorboxenvironment{topic}{
    enhanced jigsaw,
    breakable,
    colframe=TopicColor,        % 边框为靛紫色
    colback=TopicBg,           % 背景为极浅紫
    boxrule=0.5pt,             % 细边框
    leftrule=2pt,              % 左侧加粗线
    sharp corners=downhill,    % 与你其他环境一致的切角样式
    drop shadow=black!5!white, % 轻微阴影
    before skip=15pt,
    after skip=15pt
}
%----------------------------------------------------------------------------------------
%	MATH MACROS (保留你的常用定义)
%----------------------------------------------------------------------------------------
\newcommand{\g}{\mathfrak{g}}  
\newcommand{\der}{\mathrm{d}} 
\newcommand{\h}{\mathfrak{h}} 
\newcommand{\n}{\mathfrak{n}}  
\newcommand{\U}{\mathcal{U}}  
\newcommand{\C}{\mathbb{C}} 
\newcommand{\R}{\mathbb{R}} 
\newcommand{\borel}{\mathfrak{b}}
\newcommand{\Ind}{\mathrm{Ind}} 
\newcommand{\Hom}{\mathrm{Hom}}  
\newcommand{\Ext}{\mathrm{Ext}}
\newcommand{\HRule}[1]{\rule{\linewidth}{#1}}
%----------------------------------------------------------------------------------------
%	TITLE PAGE
%----------------------------------------------------------------------------------------
\begin{document}
\begin{titlepage}
    \centering
    \vspace*{2cm}
    
    \textsc{\Large Course Notes}\\[0.5cm]

    \HRule{1.5pt} \\[0.5cm]
    { \huge \bfseries Abstract Algebra III }\\[0.4cm] 
    \HRule{1.5pt} \\[1.5cm]
    
    \vspace{3cm}

    % 人员信息区域：垂直排列并区分字号
    \begin{spacing}{1.8}
        {\Large \textbf{Instructor:}} \\
        {\Large Prof. Zhicheng \textsc{Feng}} \\[0.8cm]
        
        {\normalsize \textit{Student:}} \\
        {\normalsize Leo}
    \end{spacing}

    \vfill
    {\large \today}
\end{titlepage}
\newpage
\tableofcontents
\newpage


%----------------------------------------------------------------------------------------
%    CONTENT STRUCTURE TEMPLATE
%----------------------------------------------------------------------------------------
\chapter{Lecture 1: Modules}

Message: QQ, BB. HW 30\%, Mid 20\%, Final 50\%

\section{Basic Definitions and Examples}

% --- Begin rewritten content from 抽代3.md ---
\begin{definition}
Let $R$ be a ring with $1$, and let $M$ be an abelian group. We say that $M$ is a left $R$-module if $R$ ``left acts'' on $M$:
\[
R \times M \rightarrow M, \quad (r, m) \mapsto rm
\]
such that
\[
\begin{aligned}
r(m+n) &= rm + rn, \\
(r+s)m &= rm + sm, \\
(rs)m &= r(sm), \\
1m &= m.
\end{aligned}
\quad \text{for any } r, s \in R, m,n \in M
\]

A \textbf{right $R$-module} is defined by $(r, m) \mapsto mr$ with the axiom $m(rs) = (mr)s$. If $R$ is commutative, left-modules and right-modules coincide. Otherwise, one can transform using the opposite ring $R^{\mathrm{op}}$.
\end{definition}

Consider $ M $ an $ R $-module, then $ \forall r \in R $, $ \varphi_r: M \rightarrow M $ by $ \varphi_r(m) = rm $ is an abelian group homomorphism. Let
\[
\operatorname{End}(M) = \{ \varphi: M \rightarrow M \text{ is a abelian group homomorphism } \}
\]        
Then $ \operatorname{End}(M) $ is a ring with composition as the multiplication.

\begin{theorem}
Let $M$ be an abelian group. Then $M$ is an $R$-module if and only if there is a ring homomorphism $R \rightarrow \operatorname{End}(M)$.
\end{theorem}

\begin{proof}
$(\Rightarrow)$ Suppose $M$ is an $R$-module. Define $\rho: R \rightarrow \operatorname{End}(M)$ by
\[
\rho(r)(m) = rm, \quad \forall r \in R, m \in M.
\]
For each $r \in R$, the map $\rho(r): M \rightarrow M$ is a group homomorphism since $r(m_1 + m_2) = rm_1 + rm_2$. Moreover, $\rho$ is a ring homomorphism:
\begin{itemize}
    \item $\rho(r + s)(m) = (r + s)m = rm + sm = \rho(r)(m) + \rho(s)(m)$, so $\rho(r + s) = \rho(r) + \rho(s)$.
    \item $\rho(rs)(m) = (rs)m = r(sm) = \rho(r)(\rho(s)(m))$, so $\rho(rs) = \rho(r) \circ \rho(s)$.
    \item $\rho(1)(m) = 1m = m$, so $\rho(1) = \operatorname{id}_M$.
\end{itemize}

$(\Leftarrow)$ Suppose $\rho: R \rightarrow \operatorname{End}(M)$ is a ring homomorphism. Define the scalar multiplication $R \times M \rightarrow M$ by
\[
(r, m) \mapsto \rho(r)(m).
\]
We verify the module axioms:
\begin{itemize}
    \item $r(m_1 + m_2) = \rho(r)(m_1 + m_2) = \rho(r)(m_1) + \rho(r)(m_2) = rm_1 + rm_2$.
    \item $(r + s)m = \rho(r + s)(m) = (\rho(r) + \rho(s))(m) = \rho(r)(m) + \rho(s)(m) = rm + sm$.
    \item $(rs)m = \rho(rs)(m) = (\rho(r) \circ \rho(s))(m) = \rho(r)(\rho(s)(m)) = r(sm)$.
    \item $1m = \rho(1)(m) = \operatorname{id}_M(m) = m$.
\end{itemize}
Hence $M$ is an $R$-module.
\end{proof}

\begin{example}
(1) A vector space is a module over a field.

(2) An abelian group is a module over the ring of integers $\mathbb{Z}$.
\[
\mathbb{Z} \times M \rightarrow M, \quad (n, m) \mapsto n \cdot m := m+m+\cdots+m \text{ (n times)}
\]

(3) An ideal $ I $ of a ring $ R $ is a module over $ R $.

(4) Let $M$ be an abelian group. Then $M$ is an $\operatorname{End}(M)$-module via
\[
\operatorname{End}(M) \times M \rightarrow M, \quad (\varphi, m) \mapsto \varphi(m).
\]

(5) Let $V$ be a vector space over a field $K$, and let $A: V \rightarrow V$ be a linear transformation. Then $V$ is a $K[x]$-module via
\[
K[x] \times V \rightarrow V, \quad (f(x), v) \mapsto f(A)v.
\]

(6) If $|M| = 1$, then $M$ is called the zero module.

(7) Let $ R $ be a ring with 1 and $ A \in M_n(R) $. The solution of $ AX=0 $ where $ X \in R^n $ is a right $ R $-module.   
\end{example}

\begin{definition}[$R$-submodule]
Let $R$ be a ring (with $1$) and let $M$ be a left $R$-module. A subset $N\subseteq M$
is called an \emph{$R$-submodule} of $M$ if:
\begin{enumerate}
  \item $0\in N$;
  \item $n_1,n_2\in N \implies n_1+n_2\in N$;
  \item $n\in N \implies -n\in N$;
  \item $r\in R,\ n\in N \implies rn\in N$.
\end{enumerate}
Equivalently, $N$ is a subgroup of the abelian group $(M,+)$ that is closed under the
$R$-action.
\end{definition}

\begin{example} 
(1) Let $ V $ be a vector space over a field $ K $, then any subspace of $ V $ is a $ k $-submodule of $ V $.

(2) Let $ G $ be an abelian group, then any subgroup of $ G $ is a $ \mathbb{Z} $-submodule of $ G $.

(3) Let $ R $ be a ring, $ {}_RR $ is a left (regular) $ R $-module. Then any left ideal of $ R $ is a $ R $-submodule of $ {}_RR $.

(4) Let $ G $ be an abelian group, which is an $ \operatorname{End}(G) $-module. Then any full invariant subgroup of $ G $ is a $ \operatorname{End}(G) $-submodule of $ G $.
\end{example}

\textbf{Question:} Are submodules of a finitely generated $R$-module finitely generated?

\textbf{Answer:} No, not in general. This property holds if and only if the ring $R$ is Noetherian.

\begin{example}[Submodule of a f.g. module that is not f.g.]
Let $k$ be a field and consider the polynomial ring $R = k[x_1, x_2, x_3, \dots]$ in infinitely many variables.
\begin{itemize}
    \item The ring $R$ itself can be viewed as a left $R$-module, which we denote by $M=R$. $M$ is finitely generated, for instance by the single element $1 \in R$.
    \item Consider the ideal $I = \langle x_1, x_2, x_3, \dots \rangle$ generated by all the variables. This ideal $I$ is a submodule of $M=R$.
    \item This submodule $I$ is not finitely generated. To see this, suppose for contradiction that $I$ is finitely generated by a set of polynomials $\{f_1, \dots, f_n\}$. Let $N$ be the largest index of any variable appearing in any of the polynomials $f_1, \dots, f_n$. Then any element in the ideal generated by $\{f_1, \dots, f_n\}$ must be a polynomial in variables $x_1, \dots, x_N$. However, the polynomial $x_{N+1}$ is in $I$ but cannot be expressed as a combination of $f_1, \dots, f_n$. This is a contradiction.
\end{itemize}
Therefore, the ideal $I$ is a submodule of the finitely generated $R$-module $R$ which is not itself finitely generated. This shows that the ring $R$ is not Noetherian.
\end{example}

\begin{definition} 
    Let $ M $ be an $ R $-module, $ M_i, i\in I $ are $ R $-submodules of $ M $. Then 
    \[
    \bigcap_{i\in I} M_i
    \]     
is an $ R $-submodule of $ M $.   
\[
\sum_{i\in I} = \{ \sum_{i \in I \text{ finite }} m_i | m_i \in M_i\} 
\]
is an $ R $-submodule of $ M $.
\end{definition}

\begin{definition} 
Let $M$ be an $R$-module, $ S \subseteq M $ a subset. 
\[
<S> = \{ \sum_{finite} {a_is_i} | a_i \in R, s_i \in S \}
\]
or
\[
<S>=\bigcap_{N \text{is an } R -\text{submodule of }  M,   \   S \subseteq N } N
\]
is an $ R $-submodule of $ M $, called the $ R $-submodule of $ M $ generated by $ S $.   

If $ M=<S> $, then we say $ M $ is generated by $ S $. In particular, if $ |S|=1 $ then we call $ M $ a cyclic module. If $ |S| < \infty $ then we call $ M $ a finite generated module.    
\end{definition}

\begin{definition}
    Let $R$ be a ring and let $M$ be a left $R$-module. We say that $M$ is a \emph{Noetherian module} if every ascending chain of submodules
    \[
    M_1 \subseteq M_2 \subseteq M_3 \subseteq \cdots
    \]
    stabilizes; that is, there exists $n \ge 1$ such that
    \[
    M_n = M_{n+1} = M_{n+2} = \cdots.
    \]
    Equivalently, $M$ is Noetherian if every submodule of $M$ is finitely generated.
\end{definition}



\section{Quotient Modules and Homomorphism Theorems}

\begin{definition}
    Let $R$ be a ring, let $M$ be a left $R$-module, and let $N$ be a submodule of $M$. The \emph{quotient module} (or \emph{factor module}) of $M$ by $N$ is the quotient abelian group
    \[
    M/N=\{\,x+N \mid x\in M\,\},
    \]
    with scalar multiplication defined by
    \[
    r(x+N)=rx+N
    \qquad
    (r\in R,\ x\in M).
    \]
    This is well-defined, and with this operation $M/N$ becomes a left $R$-module.
\end{definition}

\begin{example} 
(1) Let $ R $ be a ring and an $ R $-module, $ I $ be an ideal of $ R $, then $ R/I $ is a quotient $ R $-module of $ R $ by $ I $.

(2) Consider the ring $ \mathbb{Z} $ and the module $ \mathbb{Z} $ over $ \mathbb{Z} $, then for $ n \in \mathbb{Z} $, $ n \mathbb{Z} $ is a submodule of $ \mathbb{Z} $, and $ \mathbb{Z}/n\mathbb{Z} $ is a quotient $ \mathbb{Z} $-module of $ \mathbb{Z} $ by $ n\mathbb{Z} $. Furthermore, $ \mathbb{Z}/n\mathbb{Z} $ is a cyclic module.
\end{example}

\begin{definition}
    Let $R$ be a ring, and let $M$ and $N$ be left $R$-modules. A group homomorphism
    \[
    f:M\to N
    \]
    is called an \emph{$R$-module homomorphism} (or \emph{$R$-linear map}) if for all $x,y\in M$ and all $r\in R$, one has
    \[
    f(rx)=rf(x) \quad \forall r \in R, x \in M.
    \]
\end{definition}

\begin{definition} 
Let $ f $ be an $ R $-module homomorphism. The \emph{kernel} of $ f $ is the set
\[
\ker(f) = \{ x \in M \mid f(x) = 0 \}.
\]
$ \ker(f) $ is a submodule of $ M $.
The \emph{image} of $ f $ is the set
\[
\operatorname{im}(f) = \{ f(x) \in N \mid x \in M \}.
\]
$ \operatorname{im}(f) $ is a submodule of $ N $.

If $ \ker(f) = 0 $, then $ f $ is injective.
If $ \operatorname{im}(f) = N $, then $ f $ is surjective.
If $ \ker(f) = 0 $ and $ \operatorname{im}(f) = N $, then $ f $ is bijective.

A bijective $ R $-module homomorphism is called an \emph{isomorphism}.
\end{definition}

\begin{example} 
Let $ R $ be a ring and $ M,N $ be an $ R $-modules. Then $ \operatorname{Hom}_R(M,N) $ is an $ R $-module when $ R $ is commutative.     
\end{example}

\begin{theorem} 
If $ f: M\to N $ is an $ R $-module homomorphism, then we have a canonical $ R $-module isomorphism:
\[
M/ \ker(f) \rightarrow \operatorname{im}(f)
\]   
\end{theorem}

\begin{theorem} 
Let $ N $ be an $ R $-submodule of $ M $, $ \pi_N: M \to M/N $ be a quotient map. Then $ \pi_N $ induces a natural bijection:
\[
\{R-\text{submodules of } M \text{ containing } N\} \rightarrow \{R-\text{submodules of } M/N\}
\]    
\end{theorem}

\begin{theorem} 
If $ N\subseteq H$ are two $ R $-submodules of $ M $, then 
\[
M/H \cong M/N / (H/N), 
\]
\end{theorem}

\begin{theorem} 
If $ N,H $ are $ R $-submodules of $ M $, then
\[
\frac{H+N}{N} \cong \frac{H}{H \cap N} 
\] 
\end{theorem} 

\begin{proposition}
Let $M$ be an $R$-module, $N$ a submodule of $M$, $M'$ an $R$-module, and $\varphi: M \rightarrow M'$ a morphism such that $N \subseteq \ker(\varphi)$. Then there exists a unique morphism
\[
\varphi_{N}: M/N \rightarrow M'
\]
such that the following diagram commutes. Moreover, $\ker(\varphi_{N}) = \ker\varphi / N$.
\[
\begin{tikzcd}
M \arrow[r, "\varphi"] \arrow[d, "\pi_{N}"'] & M' \\
M/N \arrow[ru, "\varphi_{N}"', dashed]
\end{tikzcd}
\]
\end{proposition}

\begin{remark}
The pair $(M/N, \pi_{N})$ is unique in the following sense: if $(S, \theta)$ is a pair such that
\begin{itemize}
\item $S$ is an $R$-module,
\item $\theta: M \rightarrow S$ is a surjective morphism satisfying the property of $(M/N, \pi_{N})$ described in the above proposition,
\end{itemize}
then there exists a unique isomorphism $\psi: M/N \rightarrow S$ such that the following diagram commutes.
\[
\begin{tikzcd}[row sep=2.8em, column sep=3.8em]
M \arrow[r, two heads, "\theta"] \arrow[d, two heads, "\pi_{N}"'] & S \\
M/N \arrow[ur, dashed, "\exists!\,\psi"'] &
\end{tikzcd}
\]
\end{remark}

\begin{proof}
    Since $(S,\theta)$ satisfies the same universal property as $(M/N,\pi_N)$, we may apply that property to the canonical projection
    \[
    \pi_N : M \to M/N.
    \]
    Because $N \subseteq \ker(\pi_N)$, there exists a unique morphism
    \[
    u:S\to M/N
    \]
    such that
    \[
    u\circ \theta=\pi_N.
    \]
    
    On the other hand, since $(M/N,\pi_N)$ has the universal property of the quotient by $N$, and since $N \subseteq \ker(\theta)$, there exists a unique morphism
    \[
    \psi:M/N\to S
    \]
    such that
    \[
    \psi\circ \pi_N=\theta.
    \]
    
    We now show that $\psi$ and $u$ are inverse to each other.
    
    First, compute
    \[
    (u\circ \psi)\circ \pi_N
    = u\circ (\psi\circ \pi_N)
    = u\circ \theta
    = \pi_N.
    \]
    But the identity morphism $\operatorname{id}_{M/N}:M/N\to M/N$ also satisfies
    \[
    \operatorname{id}_{M/N}\circ \pi_N=\pi_N.
    \]
    By the uniqueness part of the universal property of $(M/N,\pi_N)$, it follows that
    \[
    u\circ \psi=\operatorname{id}_{M/N}.
    \]
    
    Similarly,
    \[
    (\psi\circ u)\circ \theta
    = \psi\circ (u\circ \theta)
    = \psi\circ \pi_N
    = \theta.
    \]
    But the identity morphism $\operatorname{id}_S:S\to S$ also satisfies
    \[
    \operatorname{id}_S\circ \theta=\theta.
    \]
    By the uniqueness part of the universal property of $(S,\theta)$, we obtain
    \[
    \psi\circ u=\operatorname{id}_S.
    \]
    
    Therefore $\psi$ is an isomorphism, with inverse $u$.
    
    Finally, the uniqueness of $\psi$ follows directly from the universal property of $(M/N,\pi_N)$: indeed, $\psi$ is the unique morphism
    \[
    \psi:M/N\to S
    \]
    such that
    \[
    \psi\circ \pi_N=\theta.
    \]
    Hence there exists a unique isomorphism
    \[
    \psi:M/N\to S
    \]
    for which the diagram commutes.
\end{proof}

\begin{proposition}
There is a unique isomorphism $\bar{\varphi}: \operatorname{coim}\varphi \rightarrow \operatorname{im}\varphi$ such that the following diagram commutes:
\begin{center}
\begin{tikzcd}
M \arrow[rr, "\varphi"] \arrow[dd, "\pi_{\operatorname{ker}(\varphi)}"'] &  & M'                                          \\
                                                                         &  &                                             \\
\operatorname{coim}(\varphi) \arrow[rr, "\overline{\varphi}"]            &  & \operatorname{im}(\varphi) \arrow[uu, hook]
\end{tikzcd}
\end{center}
Recall:
\[
\begin{aligned}
& \varphi: M \rightarrow M', \quad R\text{-morphism} \\
& \operatorname{coker}(\varphi)=M'/\operatorname{im}(\varphi) \\
& \operatorname{coim}(\varphi)=M/\ker \varphi
\end{aligned}
\]
\end{proposition}

\begin{definition}
Let $M$ be an $R$-module and $I$ an ideal of $R$. Denote by $IM$ the submodule of $M$ generated by all elements $am$ where $a \in I$, $m \in M$.
\[
IM = \left\{ \sum_{\text{finite}} a_{i} m_{i} \mid a_{i} \in I, \; m_{i} \in M \right\}
\]
\end{definition}

\begin{proposition}
Let $M$ be an $R$-module. Then $R \rightarrow \operatorname{End}(M)$ is a ring homomorphism.
\begin{enumerate}
\item $R \rightarrow \operatorname{End}(M)$ factors through the natural surjection $R \rightarrow R/I$ if and only if $IM = 0$:
\[
\begin{tikzcd}[row sep=2.8em, column sep=4.2em]
R \arrow[rr, "\rho"] \arrow[dr, two heads, "\pi_{I}"'] && \operatorname{End}(M) \\
& R/I \arrow[ur, dashed, "\exists\,\rho_{I}"'] &
\end{tikzcd}
\]
Where $\rho$ factors through $\pi_I$ means $\exists \rho_{I}: R/I \rightarrow \operatorname{End}(M)$ such that the diagram above commutes.
\item The module $M/IM$ is naturally endowed with a structure of $R/I$-module.
\item Let $\varphi: M \rightarrow M'$ be an $R$-homomorphism. Then $\varphi(IM) \subseteq IM'$, so $\varphi$ induces a morphism of $R/I$-modules
\[
M/IM \rightarrow M'/IM'.
\]
\item If $M=M_1 \oplus \cdots \oplus M_n$, then $M/IM = M_1/IM_1 \oplus \cdots \oplus M_n/IM_n$.
\item $R^n/IR^n \cong (R/IR)^n$.
\end{enumerate}
\end{proposition}

    \begin{proof}
    \textbf{(1)} Suppose first that $\rho$ factors through $\pi_I$. Then there exists a ring homomorphism
    \[
    \rho_I:R/I\to \operatorname{End}_{\mathbb{Z}}(M)
    \]
    such that
    \[
    \rho=\rho_I\circ \pi_I.
    \]
    If $a\in I$, then $\pi_I(a)=0$, so
    \[
    \rho(a)=\rho_I(0)=0.
    \]
    Therefore $am=0$ for every $m\in M$. Since $IM$ is generated by elements of the form $am$ with $a\in I$ and $m\in M$, it follows that
    \[
    IM=0.
    \]
    
    Conversely, assume that $IM=0$. Then for every $a\in I$ and every $m\in M$ we have
    \[
    \rho(a)(m)=am=0.
    \]
    Thus $\rho(a)=0$, so $I\subseteq \ker \rho$. Hence $\rho$ factors uniquely through the quotient ring $R/I$: there exists a unique ring homomorphism
    \[
    \rho_I:R/I\to \operatorname{End}_{\mathbb{Z}}(M)
    \]
    such that
    \[
    \rho=\rho_I\circ \pi_I.
    \]

    \textbf{(2)} Define an action of $R/I$ on $M/IM$ by
    \[
    (r+I)\cdot (m+IM):=rm+IM.
    \]
    We must show that this is well-defined.
    
    If $r+I=r'+I$, then $r-r'\in I$, so
    \[
    (r-r')m\in IM,
    \]
    hence
    \[
    rm+IM=r'm+IM.
    \]
    If $m+IM=m'+IM$, then $m-m'\in IM$. Since $IM$ is a submodule of $M$, we have
    \[
    r(m-m')\in IM,
    \]
    so
    \[
    rm+IM=rm'+IM.
    \]
    Thus the action is well-defined. The module axioms follow immediately from the module axioms on $M$. Therefore $M/IM$ is naturally an $R/I$-module.
    
    \textbf{(3)} Let $\varphi:M\to M'$ be an $R$-module homomorphism. Take any element of $IM$, say
    \[
    x=\sum_{j=1}^t a_j m_j
    \qquad (a_j\in I,\; m_j\in M).
    \]
    Then
    \[
    \varphi(x)=\sum_{j=1}^t \varphi(a_jm_j)=\sum_{j=1}^t a_j\varphi(m_j).
    \]
    Since each $a_j\in I$ and each $\varphi(m_j)\in M'$, it follows that
    \[
    a_j\varphi(m_j)\in IM'.
    \]
    Hence $\varphi(x)\in IM'$, and therefore
    \[
    \varphi(IM)\subseteq IM'.
    \]
    
    Thus we may define
    \[
    \overline{\varphi}:M/IM\to M'/IM',
    \qquad
    \overline{\varphi}(m+IM)=\varphi(m)+IM'.
    \]
    This is well-defined: if $m+IM=m'+IM$, then $m-m'\in IM$, so
    \[
    \varphi(m)-\varphi(m')=\varphi(m-m')\in IM',
    \]
    hence
    \[
    \varphi(m)+IM'=\varphi(m')+IM'.
    \]
    
    Finally, $\overline{\varphi}$ is $R/I$-linear, since
    \[
    \overline{\varphi}\bigl((r+I)\cdot (m+IM)\bigr)
    =\overline{\varphi}(rm+IM)
    =\varphi(rm)+IM'
    =r\varphi(m)+IM'
    =(r+I)\cdot \overline{\varphi}(m+IM).
    \]

    \textbf{(4)} Assume that
    \[
    M=M_1\oplus\cdots\oplus M_n.
    \]
    We first show that
    \[
    IM=IM_1\oplus\cdots\oplus IM_n.
    \]
    
    If $x\in IM$, then $x$ is a finite sum of elements of the form $am$ with $a\in I$ and $m\in M$. Writing
    \[
    m=m_1+\cdots+m_n
    \qquad (m_i\in M_i),
    \]
    we get
    \[
    am=am_1+\cdots+am_n,
    \]
    and each $am_i\in IM_i$. Hence $x\in IM_1+\cdots+IM_n$, so
    \[
    IM\subseteq IM_1+\cdots+IM_n.
    \]
    The reverse inclusion is immediate because each $M_i\subseteq M$, hence each $IM_i\subseteq IM$. Therefore
    \[
    IM=IM_1+\cdots+IM_n.
    \]
    
    To see that the sum is direct, suppose
    \[
    x_1+\cdots+x_n=0
    \qquad\text{with } x_i\in IM_i\subseteq M_i.
    \]
    Since
    \[
    M=M_1\oplus\cdots\oplus M_n
    \]
    is a direct sum, we must have $x_i=0$ for all $i$. Thus
    \[
    IM=IM_1\oplus\cdots\oplus IM_n.
    \]
    
    Now define
    \[
    \Psi:M_1/IM_1\oplus\cdots\oplus M_n/IM_n\to M/IM
    \]
    by
    \[
    \Psi(m_1+IM_1,\dots,m_n+IM_n)
    =(m_1+\cdots+m_n)+IM.
    \]
    This map is well-defined and $R/I$-linear. It is surjective because every element of $M$ can be written as $m_1+\cdots+m_n$. It is injective because
    \[
    (m_1+\cdots+m_n)+IM=IM
    \]
    means that
    \[
    m_1+\cdots+m_n\in IM=IM_1\oplus\cdots\oplus IM_n,
    \]
    so each $m_i\in IM_i$. Hence
    \[
    m_i+IM_i=IM_i
    \]
    for all $i$, and the kernel is trivial. Therefore $\Psi$ is an isomorphism:
    \[
    M/IM \cong M_1/IM_1 \oplus \cdots \oplus M_n/IM_n.
    \]
 
    \textbf{(5)} Since
    \[
    R^n=R\oplus\cdots\oplus R
    \]
    ($n$ copies), part (4) yields
    \[
    R^n/IR^n \cong R/IR \oplus \cdots \oplus R/IR.
    \]
    Therefore
    \[
    R^n/IR^n \cong (R/IR)^n.
    \]
    
    This completes the proof.
\end{proof}

\section{Direct Sums and Free Modules}

\begin{definition}
Direct sum and direct product $\left(\bigoplus_{i \in I} M_i \; \& \; \prod_{i \in I} M_i\right)$.

Let $ (M_i)_{i \in I} $ be a family of $ R $-modules. The direct product of $ (M_i)_{i \in I} $ is the module
\[
\prod_{i \in I} M_i = \{ \varphi: I \to \cup_{i \in I}M_i | \forall i \in I, \varphi(i) \in M_i \}
\]
with the componentwise addition and scalar multiplication. i.e.
\[
R \times \prod_{i \in I} M_i \to \prod_{i \in I} M_i, \quad (r, \varphi) \mapsto (r \varphi(i))_{i \in I}
\]
The direct sum of $ (M_i)_{i \in I} $ is the module
\[
\bigoplus_{i \in I} M_i = \{ \varphi \in \prod_{i \in I} M_i | \sharp \{ i \in I | \varphi(i) \neq 0 \} < \infty \}
\]
Equal when $|I| < \infty$. In a module category, direct product $=$ product, direct sum $=$ coproduct. 
\end{definition}

\begin{theorem} 
Suppose that $ M_1, \ldots , M_n $ are $ R $-submodules of $ M $ in which $ M = M_1 + \cdots + M_n $, $ I=[n] $. Then the followings are equivalent:
\begin{enumerate}
\item The map $ \varphi:  \sum_{i \in I} M_i \to M $ by $ \varphi( (m_i)_{i \in I} ) = \sum_{i \in I} m_i, m_i \in M_i $ is an isomorphism.
\item Any element of $ M $ can be uniquely written as a sum of elements of $ M_i $.
\item $ 0 $ can be uniquely written as a sum of elements of $ M_i $. More precisely, $ 0 = 0 + \ldots + 0 $. 
\item For $ i \in I $, $ M_i \cap (\sum_{j \neq i} M_j) = 0 $. 
\end{enumerate}
\end{theorem}

\begin{proof} 
    We prove
    \[
    (1)\Rightarrow(2)\Rightarrow(3)\Rightarrow(4)\Rightarrow(1).
    \]
    
   $(1)\Rightarrow(2)$.
    If $\varphi$ is an isomorphism, then for every $m\in M$ there exists a unique
    \[
    (m_1,\ldots,m_n)\in \sum_{i=1}^n M_i
    \]
    such that
    \[
    m=\varphi(m_1,\ldots,m_n)=m_1+\cdots+m_n.
    \]
    Hence every element of $M$ can be written uniquely as a sum of elements of the $M_i$.
    
   $(2)\Rightarrow(3)$.
    Apply $(2)$ to the element $0\in M$.
    
    $(3)\Rightarrow(4)$.
    Fix $i\in I$, and let
    \[
    x\in M_i\cap \sum_{j\ne i} M_j.
    \]
    Then $x\in M_i$, and there exist $m_j\in M_j$ for $j\ne i$ such that
    \[
    x=\sum_{j\ne i} m_j.
    \]
    Therefore
    \[
    0=x-\sum_{j\ne i} m_j
    \]
    is a decomposition of $0$ as a sum of elements from the submodules $M_1,\ldots,M_n$.
    By the uniqueness in $(3)$, this decomposition must be the trivial one. In particular,
    \[
    x=0.
    \]
    Hence
    \[
    M_i\cap \sum_{j\ne i} M_j=0.
    \]
    
  $(4)\Rightarrow(1)$.
    Since
    \[
    M=M_1+\cdots+M_n,
    \]
    the map $\varphi$ is surjective. It remains to show that $\varphi$ is injective.
    
    Suppose that
    \[
    \varphi(m_1,\ldots,m_n)=0,
    \]
    that is,
    \[
    m_1+\cdots+m_n=0
    \qquad\text{with } m_i\in M_i.
    \]
    For each $i\in I$, we have
    \[
    m_i=-\sum_{j\ne i} m_j\in \sum_{j\ne i} M_j.
    \]
    Since also $m_i\in M_i$, it follows that
    \[
    m_i\in M_i\cap \sum_{j\ne i} M_j=0.
    \]
    Thus $m_i=0$ for all $i$, so
    \[
    \ker\varphi=0.
    \]
    Therefore $\varphi$ is injective, and hence an isomorphism.
    
    This proves that the four conditions are equivalent.
\end{proof}

\begin{definition}[free module]
Let $M$ be an $R$-module.
\begin{enumerate}
\item $x_{1}, \dots, x_{n} \in M$ are linearly independent if
\[
a_{1} x_{1} + a_{2} x_{2} + \cdots + a_{n} x_{n} = 0 \Rightarrow a_{1} = a_{2} = \cdots = a_{n} = 0.
\]
\item $M$ is called free if it is generated by a linearly independent subset. This subset is called a basis of $M$.
\item If $M$ is an $R$-module with a basis $x_{1}, \ldots, x_{n}$, then $M = \langle x_{1} \rangle \oplus \cdots \oplus \langle x_{n} \rangle$.
\end{enumerate}
\end{definition}

\begin{example}
$R \oplus \cdots \oplus R = R^{n}$.
\end{example}

\begin{theorem}
Let $R$ be a commutative ring with $1$. Let $M$ be a finitely generated free $R$-module. Then any basis of $M$ has the same cardinality. The number of elements in a basis is called the rank of the free module.
\end{theorem}

\begin{proof}
Let $B=\{x_i\}_{i\in I}$ be a basis of $M$.
We first show that any basis is finite.
Since $M$ is finitely generated, choose generators $y_1,\dots,y_n$ of $M$.
Each $y_k$ is a finite $R$--linear combination of elements of $B$,
so there exists a finite subset $I_0\subset I$ such that
\[
y_1,\dots,y_n \in \sum_{i\in I_0}Rx_i.
\]
Hence
\[
M = \sum_{k=1}^n Ry_k \subseteq \sum_{i\in I_0}Rx_i \subseteq M,
\]
so $M=\sum_{i\in I_0}Rx_i$.
If $I\setminus I_0\neq\varnothing$, pick $j\in I\setminus I_0$. Then
\[
x_j \in M = \sum_{i\in I_0}Rx_i,
\]
contradicting that $B$ is linearly independent. Thus $I=I_0$ is finite.

Now let $x_1,\dots,x_r$ and $y_1,\dots,y_s$ be two bases of $M$.
Let $J\subset R$ be a maximal ideal and set $F:=R/J$, a field.
Then $M/JM$ is naturally an $F$-vector space.
Write $\overline{m}$ for the class of $m\in M$ in $M/JM$,
and $\overline{a}$ for the class of $a\in R$ in $F$.
We claim that $\overline{x_1},\dots,\overline{x_r}$ form an $F$--basis of $M/JM$.

\emph{Spanning.}
For any $m\in M$, write $m=\sum_{i=1}^r a_i x_i$. Then
\[
\overline{m} = \sum_{i=1}^r \overline{a_i}\,\overline{x_i},
\]
so $\overline{x_1},\dots,\overline{x_r}$ span $M/JM$.

\emph{Linear independence.}
Suppose
\[
\sum_{i=1}^r \overline{a_i}\,\overline{x_i} = 0 \quad \text{in } M/JM.
\]
Then $\sum_{i=1}^r a_i x_i \in JM$, so there exist $a'_1,\dots,a'_r\in J$ such that
\[
\sum_{i=1}^r a_i x_i = \sum_{i=1}^r a'_i x_i.
\]
Hence
\[
\sum_{i=1}^r (a_i - a'_i)\,x_i = 0.
\]
By linear independence of $x_1,\dots,x_r$, we get $a_i - a'_i = 0$ for all $i$,
so $a_i\in J$, thus $\overline{a_i}=0$ for all $i$.
Therefore $\overline{x_1},\dots,\overline{x_r}$ are $F$-linearly independent,
and form an $F$--basis of $M/JM$.

By the same argument, $\overline{y_1},\dots,\overline{y_s}$ also form an $F$--basis of $M/JM$. Hence
\[
r = \dim_F(M/JM) = s.
\]
Therefore any two bases of $M$ have the same cardinality.
\end{proof}

\chapter{Lecture 2}

\section{Annihilator, Torsion and Projectivity of Free Modules}

\begin{definition}
    Let $M$ be an $R$-module, $m \in M$, $Rm=<m>=\{ rm \mid r \in R \}$. Define the annihilator of $m$ in $R$ as:
\[
\operatorname{Ann}_R(m)=\{ r \in R \mid rm=0 \}.
\]
It is a left ideal of $R$.
\end{definition}

Then we have $R/\operatorname{Ann}_R(m) \cong Rm$ is an isomorphism. Which shows that all cyclic modules are isomorphic to a quotient of $R$.

\begin{definition}
    An element $m \in M$ is called
    \begin{itemize}
        \item \textbf{torsion-free} if $\operatorname{Ann}_R(m)=0$.
        \item \textbf{torsion} if $\operatorname{Ann}_R(m)\neq 0$.
    \end{itemize}
    A module is called torsion-free if all its elements except of zero are torsion-free.
\end{definition}

\begin{definition}
    The torsion part of $M$ is defined as:
    \[
    \operatorname{Tor}(M)=\{ m \in M \mid \operatorname{Ann}_R(m)\neq 0 \}.
    \]
\end{definition}

\begin{lemma}
    If $R$ is an integral domain, then the torsion submodule of $M$ is a submodule of $M$.
\end{lemma}

\begin{proof}
We verify the submodule criteria.

\emph{Non-empty.}
Since $r\cdot 0_M=0_M$ for every $r\in R$ and $R\neq 0$,
we have $\operatorname{Ann}_R(0_M)=R\neq 0$, so $0_M\in\operatorname{Tor}(M)$.

\emph{Closed under addition.}
Let $m,n\in\operatorname{Tor}(M)$.
Then there exist $\lambda,\mu\in R\setminus\{0\}$ such that $\lambda m=0$ and $\mu n=0$.
Hence
\[
(\lambda\mu)(m+n)=\mu(\lambda m)+\lambda(\mu n)=0.
\]
Since $R$ is an integral domain and $\lambda,\mu\neq 0$, we have $\lambda\mu\neq 0$,
so $\operatorname{Ann}_R(m+n)\neq 0$ and $m+n\in\operatorname{Tor}(M)$.

\emph{Closed under scalar multiplication.}
Let $m\in\operatorname{Tor}(M)$ and $r\in R$.
Choose $\lambda\in R\setminus\{0\}$ with $\lambda m=0$. Then
\[
\lambda(rm)=r(\lambda m)=0,
\]
so $\operatorname{Ann}_R(rm)\neq 0$ and $rm\in\operatorname{Tor}(M)$.

Therefore $\operatorname{Tor}(M)$ is a submodule of $M$.
\end{proof}

\begin{example}
    (1) $0 \in M$ is torsion.

    (2) Let $R$ be a commutative ring with $1$. Then $R$ is an integral domain if and only if ${}_RR$ is torsion-free. In this case $\operatorname{Tor}({}_RR)=0$, $\operatorname{Tor}(R^n)=0$ for all $n \geq 1$.

    (3) If $I$ is a nonzero ideal of $R$, then $\operatorname{Tor}(R/I)=R/I$ where $R/I$ is regarded as an $R$-module. If we regard $R/I$ as an $R/I$-module, then it is torsion-free if and only if $I$ is a prime ideal.

    (4) More generally, let $I$ be any non-zero ideal of $R$, and $M$ be an $R$-module. Then $\operatorname{Tor}(M/IM)=M/IM$ where $M/IM$ is regarded as an $R$-module. 
\end{example}

\begin{remark}
    Let $R$ be a commutative ring with $1$. Let $M$ be an $R$-module. Then $f: M \to M$ such that $f(m)=rm$ for all $m\in M$ and $r\in R$ is a homomorphism of $R$-modules. If moreover $M$ is torsion-free, then $f$ is a monomorphism.
 \end{remark}

\begin{lemma}
    If $R$ is an integral domain and $M$ is an $R$-module, then 
    \begin{enumerate}
\item $\operatorname{Tor}(\operatorname{Tor}(M)) = \operatorname{Tor}(M)$.
\item $\operatorname{Tor}(M/\operatorname{Tor}(M)) = 0$.
    \end{enumerate}
\end{lemma}

\begin{proof}\leavevmode
\begin{enumerate}
\item
Since $\operatorname{Tor}(M)$ is a submodule of $M$, its torsion submodule is
\[
\operatorname{Tor}(\operatorname{Tor}(M))
= \{\, m\in\operatorname{Tor}(M) \mid \operatorname{Ann}_R(m)\neq 0 \,\}.
\]
Every $m\in\operatorname{Tor}(M)$ already satisfies $\operatorname{Ann}_R(m)\neq 0$,
so $\operatorname{Tor}(\operatorname{Tor}(M)) \subseteq \operatorname{Tor}(M)$.
Conversely, $\operatorname{Tor}(M)\subseteq\operatorname{Tor}(\operatorname{Tor}(M))$
because the annihilator of $m$ in $R$ does not depend on the ambient module:
for any $m\in\operatorname{Tor}(M)$, its annihilator computed in the submodule
$\operatorname{Tor}(M)$ is the same as $\operatorname{Ann}_R(m)\neq 0$,
so $m\in\operatorname{Tor}(\operatorname{Tor}(M))$.
Hence $\operatorname{Tor}(\operatorname{Tor}(M))=\operatorname{Tor}(M)$.

\item
Let $\pi\colon M\to M/\operatorname{Tor}(M)$ be the canonical projection.
Write $\overline{m}=\pi(m)=m+\operatorname{Tor}(M)$.
Suppose $\overline{m}\in\operatorname{Tor}(M/\operatorname{Tor}(M))$,
i.e.\ there exists $\lambda\in R\setminus\{0\}$ such that
\[
\lambda\,\overline{m} = \overline{0} \quad\text{in } M/\operatorname{Tor}(M).
\]
This means $\lambda m\in\operatorname{Tor}(M)$,
so there exists $\mu\in R\setminus\{0\}$ with $\mu(\lambda m)=0$, i.e.
\[
(\mu\lambda)\,m = 0.
\]
Since $R$ is an integral domain and $\mu,\lambda\neq 0$, we have $\mu\lambda\neq 0$,
hence $m\in\operatorname{Tor}(M)$, so $\overline{m}=\overline{0}$.
Therefore $\operatorname{Tor}(M/\operatorname{Tor}(M))=0$.
\end{enumerate}
\end{proof}

\begin{definition}
    Let $J$ be a subscript set, denote $R^J=\prod_{j\in J} R$. $R^{(J)}=\bigoplus_{j\in J} R$.

    Let $M$ be an $R$-module. Let $x_j \in M$ where $j \in J$ be a family of elements of $M$. Then the submodule generated by $(x_j)_{j\in J}$ is denoted by $\sum_{j\in J} Rx_j$. 
\end{definition}

\begin{corollary}
    We have a morphism of $R$-modules $\phi: R^{(J)} \to M$ such that $\phi(r_j)_{j \in J}=\sum_{j \in J}r_jx_j$ for all $(r_j)_{j \in J} \in R^{(J)}$.
\begin{itemize}
\item $\phi$ is injective if and only if $\{x_j\}_{j\in J}$ are linearly independent.
\item $\phi$ is surjective if and only if $\{x_j\}_{j\in J}$ generate $M$.
\item $\phi$ is an isomorphism if and only if $M$ is free with basis $\{x_j\}_{j\in J}$. 
\end{itemize}
\end{corollary}

\begin{lemma}
    Let $L$ be a free $R$-module with basis $\{x_j\}_{j\in J}$. For any $R$-module $M$ the following map:
    \[
    \operatorname{Hom}_R(L,M) \to M^J, \quad f \mapsto (f(x_j))_{j \in J}
    \]
    is a bijection.
\end{lemma}

\begin{proposition}
    Let $L$ be a free $R$-module. 
    \begin{enumerate}
        \item Let $\pi: X\to Y$ be an epimorphism of $R$-modules. For any $R$-module homomorphism $\varphi: L \to Y$ there exists a morphism $\tilde{\varphi}: L \to X$ such that the following diagram commutes:
        \begin{center}
            \begin{tikzcd}
                & L \arrow[d, "\varphi"] \arrow[ld, "\tilde{\varphi}"'] &   \\
X \arrow[r, "\pi"'] & Y \arrow[r]                                           & 0
\end{tikzcd}
        \end{center}
        \item In particular, for any $R$-module $M$ any epimorphism $\pi: M \to L$ has a section, that is, there exists a morphism $\sigma: L \to M$ such that $\pi \sigma = \operatorname{id}_L$. In that case $\sigma$ is injective, $\operatorname{im}(\sigma) = L$ and $M = \ker \pi \oplus \operatorname{im}(\sigma)$.
        \begin{center}
            \begin{tikzcd}
                0 \arrow[r] & \operatorname{ker} \pi \arrow[r] & M \arrow[r, two heads, shift left] & L \arrow[r] \arrow[l, "\sigma", dashed, shift left] & 0
                \end{tikzcd}
        \end{center}
    \end{enumerate}
\end{proposition}

\begin{proof}
    Let $\{e_j\}_{j\in J}$ be a basis of $L$.

    \textbf{(1)} Since $\pi$ is surjective, for each $j\in J$ there exists $\tilde{x}_j \in X$ such that $\pi(\tilde{x}_j) = \varphi(e_j)$. By the lemma above, the map $\operatorname{Hom}_R(L,X) \to X^J$ given by $f \mapsto (f(e_j))_{j\in J}$ is a bijection. Hence there exists a unique $R$-module homomorphism $\tilde{\varphi}: L \to X$ such that $\tilde{\varphi}(e_j) = \tilde{x}_j$ for all $j \in J$. Then for each basis element $e_j$,
    \[
    \pi \circ \tilde{\varphi}(e_j) = \pi(\tilde{x}_j) = \varphi(e_j),
    \]
    so $\pi \circ \tilde{\varphi}$ and $\varphi$ agree on the basis $\{e_j\}_{j\in J}$. Since the basis generates $L$, we conclude $\pi \circ \tilde{\varphi} = \varphi$.

    \begin{remark}
        The uniqueness appearing in the proof of the lifting property is \emph{not} the uniqueness of the lift among all morphisms \( \tilde{\varphi}:L\to X \) satisfying
        \[
        \pi\circ \tilde{\varphi}=\varphi.
        \]
        Rather, it is only the uniqueness \emph{after a choice of lifts of the basis elements has been fixed}.
        
        Indeed, let \(L\) be a free \(R\)-module with basis \(\{e_j\}_{j\in J}\). Since \(\pi:X\to Y\) is surjective, for each \(j\in J\) we may choose an element \(\tilde{x}_j\in X\) such that
        \[
        \pi(\tilde{x}_j)=\varphi(e_j).
        \]
        By the universal property of a free module, or equivalently by the bijection
        \[
        \operatorname{Hom}_R(L,X)\xrightarrow{\sim} X^J,\qquad
        f\mapsto (f(e_j))_{j\in J},
        \]
        there exists a \emph{unique} \(R\)-module homomorphism
        \[
        \tilde{\varphi}:L\to X
        \]
        such that
        \[
        \tilde{\varphi}(e_j)=\tilde{x}_j
        \qquad \text{for all } j\in J.
        \]
        Thus the word ``unique'' here means: \emph{for the chosen family \((\tilde{x}_j)_{j\in J}\), there is a unique homomorphism realizing those values on the basis.}
        
        However, the elements \(\tilde{x}_j\) themselves are generally \emph{not} unique. If \(\tilde{x}_j\) is one lift of \(\varphi(e_j)\), then for any \(k_j\in \ker\pi\), the element
        \[
        \tilde{x}_j+k_j
        \]
        is also a lift, since
        \[
        \pi(\tilde{x}_j+k_j)=\pi(\tilde{x}_j)+\pi(k_j)=\varphi(e_j).
        \]
        Different choices of lifts of the basis elements therefore give, in general, different morphisms \( \tilde{\varphi}:L\to X \). Hence the proposition asserts only \emph{existence} of a lift, not uniqueness.
        
        Equivalently, if \(\tilde{\varphi}_1,\tilde{\varphi}_2:L\to X\) are two lifts of \(\varphi\), then
        \[
        \pi\circ \tilde{\varphi}_1=\varphi=\pi\circ \tilde{\varphi}_2,
        \]
        so
        \[
        \pi\circ (\tilde{\varphi}_1-\tilde{\varphi}_2)=0.
        \]
        Therefore
        \[
        \operatorname{im}(\tilde{\varphi}_1-\tilde{\varphi}_2)\subseteq \ker\pi,
        \]
        that is,
        \[
        \tilde{\varphi}_1-\tilde{\varphi}_2\in \operatorname{Hom}_R(L,\ker\pi).
        \]
        Conversely, if \(\delta\in \operatorname{Hom}_R(L,\ker\pi)\), then
        \[
        \pi\circ (\tilde{\varphi}_1+\delta)=\pi\circ \tilde{\varphi}_1=\varphi,
        \]
        so \(\tilde{\varphi}_1+\delta\) is again a lift of \(\varphi\). Thus the lift is unique if and only if
        \[
        \operatorname{Hom}_R(L,\ker\pi)=0.
        \]
        \end{remark}

    \textbf{(2)} Apply (1) with $X = M$, $Y = L$, and $\varphi = \operatorname{id}_L$. Then there exists $\sigma: L \to M$ such that $\pi \circ \sigma = \operatorname{id}_L$.

    \emph{$\sigma$ is injective:} If $\sigma(l) = 0$, then $l = \operatorname{id}_L(l) = \pi(\sigma(l)) = \pi(0) = 0$.

    \emph{$M = \ker\pi \oplus \operatorname{im}(\sigma)$:} For any $m \in M$, write
    \[
    m = (m - \sigma(\pi(m))) + \sigma(\pi(m)).
    \]
    The second summand $\sigma(\pi(m)) \in \operatorname{im}(\sigma)$. At first, $\pi\bigl(m - \sigma(\pi(m))\bigr) = \pi(m) - \pi\sigma(\pi(m)) = \pi(m) - \pi(m) = 0$, so it lies in $\ker\pi$. This shows $M = \ker\pi + \operatorname{im}(\sigma)$.

    For the directness, if $m \in \ker\pi \cap \operatorname{im}(\sigma)$, write $m = \sigma(l)$ for some $l \in L$. Then $0 = \pi(m) = \pi(\sigma(l)) = l$, so $m = \sigma(0) = 0$. Hence the sum is direct.
\end{proof}

\begin{exercise}
Let $L$ be a free $R$-module. $P$ a direct summand of $L$. Show that whenever $\pi: X\to Y$ is an epimorphism and $\varphi: P \to Y$ is a morphism, there exists a morphism $\tilde{\varphi}: P \to X$ such that $\varphi = \pi \circ \tilde{\varphi}$.
\end{exercise}

\section{Tensor product}

Let $R$ be a commutative ring. $M,N$ be $R$-modules. The tensor product $M \otimes_R N$ is a "linearize" of the bilinear map.

\begin{definition}
    Let $R$ be a commutative ring. $M,N,K$ be $R$-modules. $f: M\times N \to K$ is a bilinear map if for all $r\in R$, $m,m_1,m_2\in M$, $n,n_1,n_2\in N$, we have:
    \begin{itemize}
        \item $f(m_1+m_2,n) = f(m_1,n) + f(m_2,n)$
        \item $f(m,n_1+n_2) = f(m,n_1) + f(m,n_2)$
        \item $f(rm,n) = r f(m,n)$
        \item $f(m,rn) = r f(m,n)$
    \end{itemize}
\end{definition}

\begin{example}
    If $f: M\times N \to K$ is a bilinear map, then $f(x,0)=0=f(0,y)$ for all $x\in M$, $y\in N$. $f(-x,y)=-f(x,y)$ and $f(x,-y)=-f(x,y)$ for all $x\in M$, $y\in N$.
\end{example}

\begin{definition}
    The tensor product of $M$ and $N$ is a pair $(M \otimes_R N, t_{M,N})$ where $M \otimes_R N$ is an $R$-module and $t_{M,N}: M\times N \to M \otimes_R N$ is a bilinear map usually denoted by $t_{M,N}(m,n)=m \otimes n$ which satisfies the following universal properties:

    Whenever $X$ is an $R$-module and $f: M\times N \to X$ is a bilinear map, there exists a unique $R$-module homomorphism $\tilde{f}: M \otimes_R N \to X$ such that the following diagram commutes:
    \begin{center}
        \begin{tikzcd}
            M \times N \arrow[r, "f"] \arrow[d, "{t_{M,N}}"'] & X \\
            M \otimes_R N \arrow[ru, "\tilde{f}"', dashed]    &  
            \end{tikzcd}
    \end{center}
\end{definition}

\begin{theorem}
    The tensor product of $M$ and $N$ exists and is unique up to isomorphism.
\end{theorem}

\begin{proof}
    We prove the uniqueness at first.

    The tensor product is unique up to isomorphism. Suppose $(T_i, t_1)$ and $(T_2, t_2)$ are two tensor products of $M$ and $N$. By the universal property, there exists a unique $R$-module homomorphism $\varphi: T_1 \to T_2$ such that $t_2 = \varphi \circ t_1$ and a unique $R$-module homomorphism $\psi: T_2 \to T_1$ such that $t_1 = \psi \circ t_2$. Then we have:
    \[
    \varphi \circ \psi \circ t_2 = \varphi \circ t_1 = t_2 
    \]
    \[
    \psi \circ \varphi \circ t_1 = \psi \circ t_2 = t_1.
    \]
    Since such $\varphi \circ \psi$ and $\psi \circ \varphi$ are unique. It forces $\varphi \circ \psi = \operatorname{id}_{T_2}$ and $\psi \circ \varphi = \operatorname{id}_{T_1}$. Hence $\varphi$ and $\psi$ are isomorphisms. Therefore $T_1 \cong T_2$.

    Now we prove the existence.

    Let $F$ be a free $R$-module with basis $M \times N$. i.e. $F = \operatorname{span}_{R} (e_{m_i,n_i})_{m_i \in M, n_i \in N}$. Let $B$ be the submodule of $F$ generated by the elements of the following forms:
    \begin{itemize}
        \item $e_{m_1+m_2,n}-e_{m_1,n}-e_{m_2,n}$
        \item $e_{m,n_1+n_2}-e_{m,n_1}-e_{m,n_2}$
        \item $e_{rm,n}-r e_{m,n}$
        \item $e_{m,rn}-r e_{m,n}$
    \end{itemize}
    Where $m, m_1,m_2\in M$, $n,n_1,n_2\in N$, $r\in R$. Define $M\otimes_R N=F/B$, $t_{M,N}: M\times N \to M\otimes_R N$ by $t_{M,N}(m,n)=e_{m,n}+B$. It is easy to check that $t_{M,N}$ is a well-defined bilinear map. 

    We only need to show that $(M\otimes_R N, t_{M,N})$ satisfies the universal property of tensor product:

    Let $X$ be an $R$-module and $f: M\times N \to X$ a bilinear map. Since $F$ is a free $R$-module with basis $\{e_{m,n}\}_{(m,n)\in M\times N}$, by the universal property of the free module, there exists a unique $R$-module homomorphism $\hat{f}: F \to X$ such that $\hat{f}(e_{m,n}) = f(m,n)$ for all $(m,n)\in M\times N$.

    We claim that $B \subseteq \ker \hat{f}$. Indeed, since $f$ is bilinear:
    \begin{align*}
        \hat{f}(e_{m_1+m_2,n}-e_{m_1,n}-e_{m_2,n}) &= f(m_1+m_2,n)-f(m_1,n)-f(m_2,n) = 0, \\
        \hat{f}(e_{m,n_1+n_2}-e_{m,n_1}-e_{m,n_2}) &= f(m,n_1+n_2)-f(m,n_1)-f(m,n_2) = 0, \\
        \hat{f}(e_{rm,n}-r e_{m,n}) &= f(rm,n)-r f(m,n) = 0, \\
        \hat{f}(e_{m,rn}-r e_{m,n}) &= f(m,rn)-r f(m,n) = 0.
    \end{align*}
    Since $B$ is generated by elements of the above forms and $\ker\hat{f}$ is a submodule, we have $B \subseteq \ker\hat{f}$.

    By the universal property of quotient modules, $\hat{f}$ induces a unique $R$-module homomorphism $\tilde{f}: M\otimes_R N = F/B \to X$ such that $\tilde{f} \circ \pi = \hat{f}$, where $\pi: F \to F/B$ is the canonical projection. In particular, $\tilde{f}(m\otimes n) = \tilde{f}(e_{m,n}+B) = \hat{f}(e_{m,n}) = f(m,n)$, so $\tilde{f} \circ t_{M,N} = f$.

    For uniqueness, suppose $g: M\otimes_R N \to X$ is another $R$-module homomorphism with $g \circ t_{M,N} = f$. Then $g(m\otimes n) = f(m,n) = \tilde{f}(m\otimes n)$ for all $m\in M, n\in N$. Since elements of the form $m\otimes n$ generate $M\otimes_R N$ as an $R$-module, $g = \tilde{f}$.

    Therefore $(M\otimes_R N, t_{M,N})$ is a tensor product of $M$ and $N$.
\end{proof}

\begin{definition}
    The set of elements of the form $m \otimes n$ for $m \in M$ and $n \in N$ is called the set of elementary tensors. This set generates $M \otimes_R N$ as an $R$-module.
\end{definition}

\begin{example}
    Fix $m \in M$. The map $N \to M \otimes_R N, n \mapsto m \otimes n$ is a morphism of $R$-modules. Its image $ m \otimes_R N = \{m \otimes n \mid n \in N\}$ is a submodule of $M \otimes_R N$ which is isomorphic to a quotient $N/\operatorname{Ann}_R(m)N$.
\end{example}

\begin{proposition} We have the following isomorphisms: 
    \begin{enumerate}
\item $(M_1 \otimes_R M_2) \otimes_R M_3 \cong M_1 \otimes_R (M_2 \otimes_R M_3)$ where $(m_1 \otimes m_2) \otimes m_3 \mapsto m_1 \otimes (m_2 \otimes m_3)$.
\item $M_1 \otimes_R M_2 \cong M_2 \otimes_R M_1$ where $m_1 \otimes m_2 \mapsto m_2 \otimes m_1$.
\item $M_1 \otimes_R (M_2 \oplus M_3) \cong (M_1 \otimes_R M_2) \oplus (M_1 \otimes_R M_3)$ where $m_1 \otimes (m_2 \oplus m_3) \mapsto (m_1 \otimes m_2) + (m_1 \otimes m_3)$. And more generally, 
\[ 
 M_1 \otimes_R (\bigoplus_{i=1}^n M_i) \cong \bigoplus_{i=1}^n (M_1 \otimes_R M_i)
 \] 
where $m_1 \otimes (m_2 \oplus \cdots \oplus m_n) \mapsto (m_1 \otimes m_2) + \cdots + (m_1 \otimes m_n)$.
\item $R \otimes_R M \cong M$ where $r \otimes m \mapsto rm$. And more generally, $R^{(I)} \otimes_R M \cong M^{(I)}$ where $(r_i)_{i\in I} \otimes m \mapsto (r_i m)_{i\in I}$.
    \end{enumerate}
\end{proposition}

\begin{proof}
    We prove (1): $(M_1 \otimes_R M_2) \otimes_R M_3 \cong M_1 \otimes_R (M_2 \otimes_R M_3)$.

    \textbf{Step 1.} We construct $\varphi: (M_1 \otimes_R M_2) \otimes_R M_3 \to M_1 \otimes_R (M_2 \otimes_R M_3)$.

    Fix $m_3 \in M_3$. For each $m_3$, define $f_{m_3}: M_1 \times M_2 \to M_1 \otimes_R (M_2 \otimes_R M_3)$ by $f_{m_3}(m_1, m_2) = m_1 \otimes (m_2 \otimes m_3)$. This is bilinear in $(m_1, m_2)$, so by the universal property of $M_1 \otimes_R M_2$, there exists a unique $R$-module homomorphism $\tilde{f}_{m_3}: M_1 \otimes_R M_2 \to M_1 \otimes_R (M_2 \otimes_R M_3)$ such that $\tilde{f}_{m_3}(m_1 \otimes m_2) = m_1 \otimes (m_2 \otimes m_3)$:
    \begin{center}
        \begin{tikzcd}
            M_1 \times M_2 \arrow[r, "f_{m_3}"] \arrow[d, "t_{M_1,M_2}"'] & M_1 \otimes_R (M_2 \otimes_R M_3) \\
            M_1 \otimes_R M_2 \arrow[ru, "\tilde{f}_{m_3}"', dashed] &
        \end{tikzcd}
    \end{center}

    Now define $g: (M_1 \otimes_R M_2) \times M_3 \to M_1 \otimes_R (M_2 \otimes_R M_3)$ by $g(\alpha, m_3) = \tilde{f}_{m_3}(\alpha)$. We verify that $g$ is bilinear. Linearity in the first argument follows because each $\tilde{f}_{m_3}$ is an $R$-module homomorphism. For the second argument, on elementary tensors:
    \begin{align*}
        g(m_1 \otimes m_2, m_3 + m_3') &= m_1 \otimes (m_2 \otimes (m_3+m_3')) \\
        &= m_1 \otimes (m_2 \otimes m_3 + m_2 \otimes m_3') \\
        &= m_1 \otimes (m_2 \otimes m_3) + m_1 \otimes (m_2 \otimes m_3') \\
        &= g(m_1 \otimes m_2, m_3) + g(m_1 \otimes m_2, m_3'),\\
        g(m_1 \otimes m_2, rm_3) &= m_1 \otimes (m_2 \otimes rm_3) \\
        &= m_1 \otimes (r(m_2 \otimes m_3)) \\
        &= r(m_1 \otimes (m_2 \otimes m_3)) \\
        &= r\, g(m_1 \otimes m_2, m_3).
    \end{align*}
    Since elementary tensors generate $M_1 \otimes_R M_2$, bilinearity holds on all of $(M_1 \otimes_R M_2) \times M_3$.

    By the universal property of $(M_1 \otimes_R M_2) \otimes_R M_3$, there exists a unique $R$-module homomorphism
    \[
    \varphi: (M_1 \otimes_R M_2) \otimes_R M_3 \to M_1 \otimes_R (M_2 \otimes_R M_3)
    \]
    such that $\varphi((m_1 \otimes m_2) \otimes m_3) = m_1 \otimes (m_2 \otimes m_3)$:
    \begin{center}
        \begin{tikzcd}
            (M_1 \otimes_R M_2) \times M_3 \arrow[r, "g"] \arrow[d, "t"'] & M_1 \otimes_R (M_2 \otimes_R M_3) \\
            (M_1 \otimes_R M_2) \otimes_R M_3 \arrow[ru, "\varphi"', dashed] &
        \end{tikzcd}
    \end{center}

    \textbf{Step 2.} By an entirely symmetric argument (fixing $m_1$ first, then extending), we obtain an $R$-module homomorphism
    \[
    \psi: M_1 \otimes_R (M_2 \otimes_R M_3) \to (M_1 \otimes_R M_2) \otimes_R M_3
    \]
    such that $\psi(m_1 \otimes (m_2 \otimes m_3)) = (m_1 \otimes m_2) \otimes m_3$.

    \textbf{Step 3.} We show $\varphi$ and $\psi$ are inverses. On elementary tensors:
    \[
    (\psi \circ \varphi)((m_1 \otimes m_2) \otimes m_3) = \psi(m_1 \otimes (m_2 \otimes m_3)) = (m_1 \otimes m_2) \otimes m_3,
    \]
    \[
    (\varphi \circ \psi)(m_1 \otimes (m_2 \otimes m_3)) = \varphi((m_1 \otimes m_2) \otimes m_3) = m_1 \otimes (m_2 \otimes m_3).
    \]
    Since elementary tensors generate both modules, $\psi \circ \varphi = \operatorname{id}$ and $\varphi \circ \psi = \operatorname{id}$. The overall situation is summarized by the following commutative diagram:
    \begin{center}
        \begin{tikzcd}
            & M_1 \times M_2 \times M_3 \arrow[ld] \arrow[rd] & \\
            (M_1 \otimes_R M_2) \otimes_R M_3 \arrow[rr, "\varphi", shift left=1.5] & & M_1 \otimes_R (M_2 \otimes_R M_3) \arrow[ll, "\psi", shift left=1.5]
        \end{tikzcd}
    \end{center}
    where the diagonal arrows send $(m_1,m_2,m_3)$ to $(m_1\otimes m_2)\otimes m_3$ and $m_1\otimes(m_2\otimes m_3)$ respectively. Hence $\varphi$ is an isomorphism with $(m_1\otimes m_2)\otimes m_3 \mapsto m_1\otimes(m_2\otimes m_3)$.
\end{proof}

\begin{theorem}
    Let $M, N, P$ be $R$-modules. We have a canonical isomorphism:
    \[
    \phi_{M,N,P}: \operatorname{Hom}_R(M \otimes_R N, P) \to \operatorname{Hom}_R(M, \operatorname{Hom}_R(N, P))
    \]
    such that $\phi_{M,N,P}(f)(m)(n) = f(m \otimes n)$.
Furthermore, if $R$ is commutative, then $\phi_{M,N,P}$ is an $R$-module isomorphism. If $R$ is not commutative, then $\phi_{M,N,P}$ is only a $\mathbb{Z}$-module isomorphism.
\end{theorem}

\begin{proof}
    Let $f \in \operatorname{Hom}_R(M \otimes_R N, P)$. For each $m \in M$, define $\phi_{M,N,P}(f)(m): N \to P$ by
    \[
    \phi_{M,N,P}(f)(m)(n) = f(m \otimes n).
    \]
    We verify that $\phi_{M,N,P}(f)(m) \in \operatorname{Hom}_R(N, P)$: for $n_1, n_2 \in N$ and $r \in R$,
    \begin{align*}
        \phi_{M,N,P}(f)(m)(n_1 + n_2) &= f(m \otimes (n_1+n_2)) = f(m \otimes n_1 + m \otimes n_2) \\
        &= f(m\otimes n_1) + f(m\otimes n_2),\\
        \phi_{M,N,P}(f)(m)(rn) &= f(m \otimes rn) = f(r(m\otimes n)) \\
        &= rf(m\otimes n) = r\,\phi_{M,N,P}(f)(m)(n).
    \end{align*}
    Next, we verify that $\phi_{M,N,P}(f) \in \operatorname{Hom}_R(M, \operatorname{Hom}_R(N,P))$: for $m_1, m_2 \in M$, $r \in R$, and any $n\in N$,
    \begin{align*}
        \phi_{M,N,P}(f)(m_1+m_2)(n) &= f((m_1+m_2)\otimes n) \\
        &= f(m_1\otimes n) + f(m_2\otimes n)\\
        \phi_{M,N,P}(f)(rm)(n) &= f(rm\otimes n) \\
        &= f(r(m\otimes n)) \\
        &= rf(m\otimes n) \\
        &= (r\,\phi_{M,N,P}(f)(m))(n).
    \end{align*}
    So 
    \[
        \phi_{M,N,P}(f)(m_1+m_2) = \phi_{M,N,P}(f)(m_1)+\phi_{M,N,P}(f)(m_2)
   \]
        and $\phi_{M,N,P}(f)(rm) = r\,\phi_{M,N,P}(f)(m)$, confirming $R$-linearity.

     Let $g \in \operatorname{Hom}_R(M, \operatorname{Hom}_R(N, P))$. Define $\hat{g}: M \times N \to P$ by $\hat{g}(m,n) = g(m)(n)$. We check that $\hat{g}$ is $R$-bilinear:
    \begin{align*}
        \hat{g}(m_1+m_2,n) &= g(m_1+m_2)(n) = g(m_1)(n)+g(m_2)(n) = \hat{g}(m_1,n)+\hat{g}(m_2,n),\\
        \hat{g}(m,n_1+n_2) &= g(m)(n_1+n_2) = g(m)(n_1)+g(m)(n_2) = \hat{g}(m,n_1)+\hat{g}(m,n_2),\\
        \hat{g}(rm,n) &= g(rm)(n) = (rg(m))(n) = rg(m)(n) = r\,\hat{g}(m,n),\\
        \hat{g}(m,rn) &= g(m)(rn) = rg(m)(n) = r\,\hat{g}(m,n).
    \end{align*}
    By the universal property of $M \otimes_R N$, there exists a unique $R$-module homomorphism $\psi_{M,N,P}(g): M\otimes_R N \to P$ such that $\psi_{M,N,P}(g)(m\otimes n) = g(m)(n)$:
    \begin{center}
        \begin{tikzcd}
            M \times N \arrow[r, "\hat{g}"] \arrow[d, "t_{M,N}"'] & P \\
            M \otimes_R N \arrow[ru, "{\psi_{M,N,P}(g)}"', dashed] &
        \end{tikzcd}
    \end{center}

    For any $f \in \operatorname{Hom}_R(M\otimes_R N,P)$:
    \[
    \psi_{M,N,P}(\phi_{M,N,P}(f))(m\otimes n) = \phi_{M,N,P}(f)(m)(n) = f(m\otimes n).
    \]
    Since elementary tensors generate $M\otimes_R N$, we have $\psi_{M,N,P}(\phi_{M,N,P}(f)) = f$.

    For any $g \in \operatorname{Hom}_R(M, \operatorname{Hom}_R(N,P))$, for all $m\in M$ and $n\in N$:
    \[
    \phi_{M,N,P}(\psi_{M,N,P}(g))(m)(n) = \psi_{M,N,P}(g)(m\otimes n) = g(m)(n).
    \]
    So $\phi_{M,N,P}(\psi_{M,N,P}(g))(m) = g(m)$ for all $m$, hence $\phi_{M,N,P}(\psi_{M,N,P}(g)) = g$.

    Therefore $\phi_{M,N,P}$ is a bijection of sets (in fact, of abelian groups).

    It remains to show that $\phi_{M,N,P}$ respects the appropriate module structure.

    Both $\operatorname{Hom}_R(M\otimes_R N, P)$ and $\operatorname{Hom}_R(M, \operatorname{Hom}_R(N,P))$ are abelian groups under pointwise addition. We verify $\phi_{M,N,P}$ is a group homomorphism: for $f_1, f_2 \in \operatorname{Hom}_R(M\otimes_R N, P)$,
    \begin{align*}
    \phi_{M,N,P}(f_1+f_2)(m)(n) &= (f_1+f_2)(m\otimes n) \\
    &= f_1(m\otimes n)+f_2(m\otimes n) \\
    &= \phi_{M,N,P}(f_1)(m)(n)+\phi_{M,N,P}(f_2)(m)(n)
    \end{align*}
    so $\phi_{M,N,P}(f_1+f_2)=\phi_{M,N,P}(f_1)+\phi_{M,N,P}(f_2)$. Hence $\phi_{M,N,P}$ is a $\mathbb{Z}$-module isomorphism.

    If $R$ is commutative, $\operatorname{Hom}_R(M\otimes_R N, P)$ carries an $R$-module structure via $(r\cdot f)(x) = rf(x)$, and similarly for $\operatorname{Hom}_R(M, \operatorname{Hom}_R(N,P))$. Then
    \begin{align*}
    \phi_{M,N,P}(rf)(m)(n) &= (rf)(m\otimes n) \\
    &= r f(m\otimes n) \\
    &= r\,\phi_{M,N,P}(f)(m)(n) \\
    &= (r\,\phi_{M,N,P}(f))(m)(n)
    \end{align*}
    so $\phi_{M,N,P}(rf) = r\,\phi_{M,N,P}(f)$, and $\phi_{M,N,P}$ is an $R$-module isomorphism.

    If $R$ is not commutative, $\operatorname{Hom}_R(M\otimes_R N, P)$ does not naturally carry an $R$-module structure (scalar multiplication by $r$ need not commute with $R$-linearity), so $\phi_{M,N,P}$ is only a $\mathbb{Z}$-module isomorphism.
\end{proof}

\begin{definition}
    Let $f: M\to M'$ and $g: N\to N'$ be $R$-module homomorphisms. The tensor product of $f$ and $g$ is the $R$-module homomorphism:
    \[
    f \otimes g: M \otimes_R N \to M' \otimes_R N'
    \]
    defined by $(f \otimes g)(m \otimes n) = f(m) \otimes g(n)$.
\end{definition}

\begin{definition}
    Let $M_i, i \in I$ be a family of $R$-modules. Use the following notation:
    \[
    L^n(M_1, \ldots, M_n; N) = \{ f: M_1 \times \cdots \times M_n \to N \mid f \text{ is $n$-linear } \}
    \]
    In particular,
    \[
    L^2(M_1, M_2; N) \cong L(M_1 \otimes M_2, N) \cong \operatorname{Hom}_R(M_1 \otimes M_2, N)
    \]
    or equivalently,
    \[
    L^2(M_1 \otimes M_2, N) \cong L(M_1, L(M_2, N))
    \]
\end{definition}

\begin{proposition}
    Let $M,N$ be two finitely generated free $R$-modules. Denote 
    \[
    M^{\vee} = \operatorname{Hom}_R(M, R)
    \]
    Then we have a canonical isomorphism:
    \[
    M^{\vee} \otimes N \cong \operatorname{Hom}_R(M, N)
    \]
\end{proposition}

\begin{proof}
    Since $M$ is a finitely generated free $R$-module, we have $M \cong R^m$ for some $m \geq 1$. We construct an explicit isomorphism
    \[
    \Phi: M^{\vee} \otimes_R N \to \operatorname{Hom}_R(M, N).
    \]
    Define a bilinear map $\beta: M^{\vee} \times N \to \operatorname{Hom}_R(M,N)$ by
    \[
    \beta(f, n)(m) = f(m)\, n, \qquad f \in M^{\vee},\; n \in N,\; m \in M.
    \]
    We check $R$-bilinearity. For $f_1, f_2 \in M^{\vee}$, $n, n_1, n_2 \in N$, $m \in M$, and $r \in R$:
    \begin{align*}
        \beta(f_1+f_2, n)(m) &= (f_1+f_2)(m)\,n = f_1(m)\,n + f_2(m)\,n \\
        &= \beta(f_1,n)(m)+\beta(f_2,n)(m),\\
        \beta(f, n_1+n_2)(m) &= f(m)(n_1+n_2) = f(m)\,n_1+f(m)\,n_2 \\
        &= \beta(f,n_1)(m)+\beta(f,n_2)(m),\\
        \beta(rf, n)(m) &= (rf)(m)\,n = r\,f(m)\,n = r\,\beta(f,n)(m),\\
        \beta(f, rn)(m) &= f(m)\,(rn) = r\,f(m)\,n = r\,\beta(f,n)(m).
    \end{align*}
    By the universal property of the tensor product, there exists a unique $R$-module homomorphism
    \[
    \Phi: M^{\vee} \otimes_R N \to \operatorname{Hom}_R(M, N), \qquad \Phi(f \otimes n)(m) = f(m)\,n.
    \]

    Now let $\{e_1, \ldots, e_m\}$ be a basis of $M \cong R^m$, and let $\{e_1^*, \ldots, e_m^*\}$ be the dual basis, i.e.\ $e_i^* \in M^{\vee}$ with $e_i^*(e_j) = \delta_{ij}$. Define
    \[
    \Psi: \operatorname{Hom}_R(M, N) \to M^{\vee} \otimes_R N, \qquad \Psi(g) = \sum_{i=1}^{m} e_i^* \otimes g(e_i).
    \]
    We verify $\Psi$ is an $R$-module homomorphism: for $g_1, g_2 \in \operatorname{Hom}_R(M,N)$ and $r \in R$,
    \begin{align*}
        \Psi(g_1+g_2) &= \sum_{i=1}^m e_i^* \otimes (g_1+g_2)(e_i) = \sum_{i=1}^m e_i^* \otimes (g_1(e_i)+g_2(e_i)) \\
        &= \sum_{i=1}^m e_i^* \otimes g_1(e_i) + \sum_{i=1}^m e_i^* \otimes g_2(e_i) = \Psi(g_1)+\Psi(g_2),\\
        \Psi(rg) &= \sum_{i=1}^m e_i^* \otimes (rg)(e_i) = \sum_{i=1}^m e_i^* \otimes r\,g(e_i) \\
        &= r \sum_{i=1}^m e_i^* \otimes g(e_i) = r\,\Psi(g).
    \end{align*}

    We show $\Phi$ and $\Psi$ are mutual inverses. For any $g \in \operatorname{Hom}_R(M,N)$ and $m = \sum_j r_j e_j \in M$:
    \begin{align*}
        \Phi(\Psi(g))(m) &= \Phi\!\Bigl(\sum_{i=1}^m e_i^* \otimes g(e_i)\Bigr)\!\Bigl(\sum_j r_j e_j\Bigr) \\
        &= \sum_{i=1}^m e_i^*\!\Bigl(\sum_j r_j e_j\Bigr) g(e_i) = \sum_{i=1}^m r_i\, g(e_i) \\
        &= g\!\Bigl(\sum_{i=1}^m r_i e_i\Bigr) = g(m).
    \end{align*}
    So $\Phi \circ \Psi = \operatorname{id}$.

    For any elementary tensor $f \otimes n \in M^{\vee} \otimes_R N$:
    \begin{align*}
        \Psi(\Phi(f \otimes n)) &= \sum_{i=1}^m e_i^* \otimes \Phi(f \otimes n)(e_i) = \sum_{i=1}^m e_i^* \otimes f(e_i)\,n \\
        &= \sum_{i=1}^m f(e_i)\, (e_i^* \otimes n) = \Bigl(\sum_{i=1}^m f(e_i)\, e_i^*\Bigr) \otimes n = f \otimes n,
    \end{align*}
    where the last equality uses $f = \sum_{i=1}^m f(e_i)\, e_i^*$ (expansion of $f$ in the dual basis). Since elementary tensors generate $M^{\vee} \otimes_R N$ and $\Psi \circ \Phi$ is $R$-linear, we have $\Psi \circ \Phi = \operatorname{id}$.

    Therefore $\Phi$ is an $R$-module isomorphism, giving $M^{\vee} \otimes_R N \cong \operatorname{Hom}_R(M, N)$.
\end{proof}

\begin{proposition}
    Assume that $R$ is a subring of $S$, and $M$ is an $R$-module. Then the bilinear map:
    \[
    S \times (S \otimes_R M) \to S \otimes_R M
    \]
    defined by
    \[
    (s, s' \otimes m) \mapsto s s' \otimes m
    \]
    defines a structure of $S$-module on the $R$-module $S \otimes_R M$.

    More generally, let $f: R \to T$ be a ring homomorphism. And view $T$ as an $R$-module via $R x T \to T, (\lambda, \mu) \mapsto f(\lambda) \mu$. Then the bilinear map:
    \[
    T \times (T \otimes_R M) \to T \otimes_R M
    \]
    defined by
    \[
    (t, t' \otimes m) \mapsto t t' \otimes m
    \]
    defines a structure of $T$-module on the $R$-module $T \otimes_R M$.
 \end{proposition}

\begin{proposition}
    Let $I$ be an ideal of $R$. Let $M$ be an $R$-module. Then the morphism of $R$-modules:
    \[
    M \to (R/I) \otimes_R M
    \]
    defined by
    \[
    m \mapsto 1_{R/I} \otimes m
    \]
  induces an isomorphism of $R$-modules:
  \[
  M/(IM) \cong (R/I) \otimes_R M
  \]
  as $R/I$-modules.
\end{proposition}

\section{Balanced product}

\begin{definition}
    Let $R$ be a ring (not necessarily commutative). Let $M_R$ be a right $R$-module and $_{R}N$ be a left $R$-module. The balanced product of $M$ and $N$ is defined as a pair $(P,f)$ where $P$ is an abelian group and $f: M \times N \to P$ is a bilinear map satisfies:
    \begin{itemize}
\item $f(rm,n) = f(m,rn)$ for all $r \in R$, $m \in M$, $n \in N$.
\item $f(m+m',n) = f(m,n) + f(m',n)$ for all $m,m' \in M$, $n \in N$.
\item $f(m,n+n') = f(m,n) + f(m,n')$ for all $m \in M$, $n,n' \in N$.
    \end{itemize}
\end{definition}


\begin{definition}
    A tensor product of $M_R$ and $_{R}N$ is a balanced product $(M \otimes_R N,t_{M,N})$ such that if $(P,f)$ is any balanced product of $M_R$ and $_{R}N$, then there exists a unique abelian group homomorphism $\tilde{f}: M \otimes_R N \to P$ such that the following diagram commutes:
    \begin{center}
        \begin{tikzcd}
            M \times N \arrow[r, "{t_{M,N}}"] \arrow[d, "f"'] & M \otimes_R N \arrow[ld, "\tilde{f}"', dashed] \\
            P                                                 &                                               
            \end{tikzcd}
    \end{center}
\end{definition}






\chapter{Homework}

\section{Homework Week 1}

\begin{exercise}
    Let $R$ be a ring, and let $M$ be an $R$-module. Suppose that $S$ is a ring and
    \[
        \theta : S \to R
    \]
    is a ring homomorphism satisfying $\theta(1_S)=1_R$. Let $S$ act on $M$ by
    \[
        sx=\theta(s)x,
    \]
    where $x\in M$ and $s\in S$. Prove that $M$ is an $S$-module.
\end{exercise}

\begin{proof}
    This action is well-defined because $\forall s \in S$, $\theta(s) \in R$ and
    $\theta(s)x \in M$ for all $x \in M$. We need to verify the module axioms, for all
    $s, t \in S$ and $x, y \in M$:
    \begin{itemize}
        \item $s(x+y) = \theta(s)(x+y) = \theta(s)x + \theta(s)y = sx + sy$
        \item $(s+t)x = \theta(s+t)x = \theta(s)x + \theta(t)x = sx + tx$
        \item $(st)x=\theta(st)x=\theta(s)\theta(t)x=\theta(s)(\theta(t)x)=s(tx)$.
        \item $1_S x = \theta(1_S)x = 1_R x = x$
    \end{itemize}
    Hence $M$ is an $S$-module.
\end{proof}

\begin{exercise}
    Let $R$ be a ring, and let $M$ and $N$ be two right $R$-modules. Let
    $\operatorname{Hom}_{\mathbb{Z}}(M,N)$ be the set of group homomorphisms from $M$ to
    $N$ (as abelian groups). The action of $R$ on
    $\operatorname{Hom}_{\mathbb{Z}}(M,N)$ is given as follows: for any $a\in R$ and
    $\varphi\in \operatorname{Hom}_{\mathbb{Z}}(M,N)$, define
    \[
        (\varphi a)(x)=\varphi(xa), \qquad \forall x\in M.
    \]
    Prove that $\operatorname{Hom}_{\mathbb{Z}}(M,N)$ is a left $R$-module.
\end{exercise}

\begin{proof}
    This action is well-defined. $\forall r \in R$, $\varphi r \in
    \operatorname{Hom}_{\mathbb{Z}}(M,N)$ for all $\varphi \in
    \operatorname{Hom}_{\mathbb{Z}}(M,N)$ as $(\varphi r)(x) = \varphi(xr) \in N$ for all
    $x \in M$. We need to verify the module axioms, for all $r_1, r_2 \in R$ and
    $\varphi, \psi \in \operatorname{Hom}_{\mathbb{Z}}(M,N)$, $\forall x \in M$:
    \begin{itemize}
        \item $((\varphi + \psi)r)(x) = (\varphi + \psi)(xr) = \varphi(xr) + \psi(xr) = (\varphi r)(x) + (\psi r)(x)$
        \item $(\varphi(r_1 + r_2))(x) = \varphi(x(r_1 + r_2)) = \varphi(xr_1 + xr_2) = \varphi(xr_1) + \varphi(xr_2) = (\varphi r_1)(x) + (\varphi r_2)(x)$
        \item $((\varphi r_1)r_2)(x) = (\varphi r_2)(xr_1) = \varphi(xr_2r_1) = (\varphi r_2 r_1)(x)$
        \item $( \varphi 1_R)(x) =  \varphi(x 1_R) = \varphi(x) $
    \end{itemize}
    Hence $\operatorname{Hom}_{\mathbb{Z}}(M,N)$ is a left $R$-module.
\end{proof}

\begin{exercise}
    Let $R$ be a ring, and let $M$ be an $R$-module. Denote
    \[
        \operatorname{Ann}_R M=\{r\in R \mid rm=0,\ \forall m\in M\}.
    \]
    We call $M$ a faithful $R$-module if $\operatorname{Ann}_R M=\{0\}$. Define the action
    of the quotient ring $R/\operatorname{Ann}_R M$ on $M$ by
    \[
        \overline{r}\,x=rx,
    \]
    where $x\in M$, $r\in R$, and $\overline{r}=r+\operatorname{Ann}_R M$. Prove that $M$
    is a faithful $R/\operatorname{Ann}_R M$-module.
\end{exercise}

\begin{proof}
    First we show that the action is well-defined. Suppose
    \[
        \bar r=\bar r' \in R/\operatorname{Ann}_R M.
    \]
    Then $r-r'\in \operatorname{Ann}_R M$, so for every $x\in M$,
    \[
        (r-r')x=0.
    \]
    Hence $rx=r'x$, so $\bar r\,x$ does not depend on the choice of representative.

    The module axioms follow immediately from those of the $R$-module structure on $M$.

    Now we prove faithfulness. Suppose $\bar r\in R/\operatorname{Ann}_R M$ acts trivially on
    $M$, that is,
    \[
        \bar r\,x=0 \qquad \forall x\in M.
    \]
    By definition, this means
    \[
        rx=0 \qquad \forall x\in M.
    \]
    Therefore $r\in \operatorname{Ann}_R M$, so $\bar r=\bar 0$ in the quotient ring. Hence
    \[
        \operatorname{Ann}_{R/\operatorname{Ann}_R M}(M)=\{\bar 0\},
    \]
    and $M$ is faithful as an $R/\operatorname{Ann}_R M$-module.
\end{proof}

\begin{exercise}
    Let $\varphi : M \to M$ be a homomorphism of $R$-modules such that
    \[
        \varphi\varphi=\varphi.
    \]
    Prove that
    \[
        M=\ker\varphi \oplus \operatorname{im}\varphi.
    \]
\end{exercise}

\begin{proof}
    Let $x \in M$. Then $x = \varphi(x) + (x - \varphi(x))$. Where $\varphi(x) \in
    \operatorname{im}\varphi$ and $x - \varphi(x) \in \ker\varphi$. Hence $M =
    \operatorname{im}\varphi + \ker\varphi$. We only need to show that $\ker\varphi \cap
    \operatorname{im}\varphi = \{0\}$.
    Suppose $x \in \ker\varphi \cap \operatorname{im}\varphi$. Then $x = \varphi(y)$ for some
    $y \in M$ and $\varphi(x) = 0$. Then $\varphi(x) = \varphi \varphi(y) = \varphi(y)=x=0$,
    so $x = 0$. Hence $\ker\varphi \cap \operatorname{im}\varphi = \{0\}$.
\end{proof}

\begin{exercise}
    Show that any $R$-module is a homomorphism image of some free $R$-module.
\end{exercise}

\begin{proof}
    Let $M$ be a left $R$-module. Consider the free $R$-module
    \[
        F=\bigoplus_{x\in M} Re_x
    \]
    with basis $\{e_x\}_{x\in M}$ indexed by the underlying set of $M$.
    Define an $R$-module homomorphism
    \[
        \pi:F\to M
    \]
    by
    \[
        \pi(e_x)=x \qquad (x\in M),
    \]
    and extend $R$-linearly. Thus
    \[
        \pi\Bigl(\sum_{i=1}^n r_i e_{x_i}\Bigr)=\sum_{i=1}^n r_i x_i.
    \]
    This is well-defined and $R$-linear.

    For any $m\in M$, we have
    \[
        \pi(e_m)=m,
    \]
    so $\pi$ is surjective. Therefore $M$ is a homomorphic image of the free module $F$.
    Equivalently,
    \[
        M\cong F/\ker\pi.
    \]
\end{proof}

\begin{exercise}
    Let $\mathbb{Q}$ be the field of rational numbers. Is the additive group
    $(\mathbb{Q},+)$ a free $\mathbb{Z}$-module?
\end{exercise}

\begin{proof}
    Claim: any two non-zero element of $\mathbb{Q}$ are linearly dependent over
    $\mathbb{Z}$. Hence if $\mathbb{Q}$ is a free $\mathbb{Z}$-module, then it must have a
    basis of cardinality 1.

    Proof of claim: let $0 \neq x = \frac{a}{b}$ and $0 \neq y =\frac{c}{d}$ be two
    elements of $\mathbb{Q}$. Then $dcx - aby=0$. If $a=0$ or $c=0$, then $x=0$ or $y=0$,
    which contradicts the assumption. Hence $a, c \neq 0$ which means $dc,ab$ are two
    non-zero integer coefficients. Hence $x$ and $y$ are linearly dependent over
    $\mathbb{Z}$. And if $\mathbb{Q}$ is a free $\mathbb{Z}$-module, then it must have a
    basis of cardinality 1.

    But if $\mathbb{Q}$ is a cyclic $\mathbb{Z}$-module, then it has only one generator.
    Suppose it is $z=\frac{e}{f}$. Then $\mathbb{Q}=\{n \frac{e}{f} \mid n \in \mathbb{Z}\}$.
    But this is not possible because we can easily pick $\frac{1}{2f} \notin \mathbb{Q}$,
    which is absurd.

    Therefore $\mathbb{Q}$ is not a free $\mathbb{Z}$-module.
\end{proof}

\section{Homework Week 2}

\begin{exercise}
    Let $R$ be a commutative ring, and suppose that $\varphi$ is an endomorphism of the free $R$-module $R^m$. Prove that if $\varphi$ is surjective, then $\varphi$ is injective. Consider that: if $\varphi$ is injective, is $\varphi$ surjective?
\end{exercise}

\begin{proof}
    Let $\{e_1,\ldots,e_m\}$ be the standard basis of $R^m$, and let $A\in M_m(R)$ be the matrix of $\varphi$ with respect to this basis. Since $\varphi$ is surjective, for each $j$ there exists $v_j\in R^m$ with $\varphi(v_j)=e_j$. Write $v_j=\sum_{i=1}^m b_{ij}e_i$ and let $B=(b_{ij})\in M_m(R)$. Then $AB=I_m$, so
    \[
        \det(A)\det(B)=\det(AB)=1.
    \]
    Hence $\det(A)$ is a unit in $R$. It follows that $A$ is invertible (with $A^{-1}=\det(A)^{-1}\operatorname{adj}(A)$), so $\varphi$ is an isomorphism. In particular $\varphi$ is injective.

    \medskip
    \noindent\textbf{Counterexample for the converse.} The converse is false in general. Let $R = \mathbb{Z}$ and $m = 1$, so $R^m = \mathbb{Z}$. Define $\varphi: \mathbb{Z} \to \mathbb{Z}$ by $\varphi(n) = 2n$. Then $\varphi$ is injective but not surjective (e.g.\ $1 \notin \operatorname{im}\varphi$).
\end{proof}

\begin{exercise}
    Let $\varphi: M \to N$ be a homomorphism of $R$-modules. Prove that:
    \begin{enumerate}
        \item $\varphi$ is a monomorphism if and only if for any $R$-module $T$ and any two $R$-module homomorphisms $\xi, \eta: T \to M$, $\varphi\xi = \varphi\eta$ implies $\xi = \eta$;
        \item $\varphi$ is an epimorphism if and only if for any $R$-module $S$ and any two $R$-module homomorphisms $\rho, \theta: N \to S$, $\rho\varphi = \theta\varphi$ implies $\rho = \theta$.
    \end{enumerate}
\end{exercise}

\begin{proof}\leavevmode
    \begin{enumerate}
        \item ($\Rightarrow$) Suppose $\varphi$ is injective and $\varphi\xi=\varphi\eta$ for $\xi,\eta:T\to M$. Then for every $t\in T$, $\varphi(\xi(t))=\varphi(\eta(t))$, so $\xi(t)-\eta(t)\in\ker\varphi=\{0\}$. Hence $\xi(t)=\eta(t)$ for all $t$, i.e.\ $\xi=\eta$.

        ($\Leftarrow$) We show $\ker\varphi=\{0\}$. Let $T=\ker\varphi$, $\xi:T\hookrightarrow M$ be the inclusion, and $\eta=0:T\to M$ the zero map. Then $\varphi\xi(t)=\varphi(t)=0=\varphi\eta(t)$ for all $t\in T$, so $\varphi\xi=\varphi\eta$. By hypothesis $\xi=\eta$, which means $t=\xi(t)=\eta(t)=0$ for all $t\in T$. Hence $\ker\varphi=\{0\}$, so $\varphi$ is injective.

        \item ($\Rightarrow$) Suppose $\varphi$ is surjective and $\rho\varphi=\theta\varphi$. For any $n\in N$, pick $m\in M$ with $\varphi(m)=n$. Then $\rho(n)=\rho(\varphi(m))=\theta(\varphi(m))=\theta(n)$. Hence $\rho=\theta$.

        ($\Leftarrow$) We show $\varphi$ is surjective. Let $S=N/\operatorname{im}\varphi$, $\rho:N\to S$ the canonical projection, and $\theta=0:N\to S$. For any $m\in M$, $\rho(\varphi(m))=\overline{\varphi(m)}=\bar 0=\theta(\varphi(m))$, so $\rho\varphi=\theta\varphi$. By hypothesis $\rho=\theta$, i.e.\ the projection $N\to N/\operatorname{im}\varphi$ is the zero map. This means $N=\operatorname{im}\varphi$, so $\varphi$ is surjective.
    \end{enumerate}
\end{proof}

\begin{exercise}
    Let $M$ and $M_i\ (1 \le i \le n)$ be $R$-modules. Show that $M \simeq \bigoplus_{i=1}^{n} M_i$ if and only if for every $i \in \{1, \ldots, n\}$, there exists an $\varphi_i \in \operatorname{End}_R(M)$ such that $\operatorname{im}\varphi_i \simeq M_i$, $\varphi_i\varphi_j = 0$ for $i \neq j$, $\varphi_i^2 = \varphi_i$, and $\varphi_1 + \varphi_2 + \cdots + \varphi_n = \operatorname{id}_M$.
\end{exercise}

\begin{proof}
    ($\Rightarrow$) Suppose $f:M\xrightarrow{\sim}\bigoplus_{i=1}^n M_i$. Let $\iota_i:M_i\hookrightarrow \bigoplus_{j=1}^n M_j$ be the canonical inclusion and $\pi_i:\bigoplus_{j=1}^n M_j\twoheadrightarrow M_i$ the canonical projection. Define
    \[
        \varphi_i = f^{-1}\circ\iota_i\circ\pi_i\circ f \in \operatorname{End}_R(M).
    \]
    Then:
    \begin{itemize}
        \item $\varphi_i^2 = f^{-1}\iota_i\pi_i f f^{-1}\iota_i\pi_i f = f^{-1}\iota_i(\pi_i\iota_i)\pi_i f = f^{-1}\iota_i\pi_i f = \varphi_i$, since $\pi_i\iota_i=\operatorname{id}_{M_i}$.
        \item For $i\neq j$: $\varphi_i\varphi_j = f^{-1}\iota_i(\pi_i\iota_j)\pi_j f = 0$, since $\pi_i\iota_j=0$ when $i\neq j$.
        \item $\sum_{i=1}^n \varphi_i = f^{-1}\left(\sum_{i=1}^n \iota_i\pi_i\right) f = f^{-1}\circ\operatorname{id}\circ f = \operatorname{id}_M$, since $\sum_i \iota_i\pi_i = \operatorname{id}_{\bigoplus M_j}$.
        \item $\operatorname{im}\varphi_i = f^{-1}\circ\iota_i\circ\pi_i\circ f(M) \simeq M_i$.
    \end{itemize}

    ($\Leftarrow$) Suppose such $\varphi_i$ exist. Define $g:M\to \bigoplus_{i=1}^n \operatorname{im}\varphi_i$ by $g(x)=(\varphi_1(x),\ldots,\varphi_n(x))$.

    \textit{Surjective}: for any $(y_1,\ldots,y_n)\in\bigoplus \operatorname{im}\varphi_i$, let $x = y_1+\cdots+y_n \in M$. Then
    \[
        \varphi_j(x) = \sum_{i=1}^n \varphi_j(y_i) = \sum_{i=1}^n \varphi_j(\varphi_i(x_i)) = \varphi_j(\varphi_j(x_j)) = \varphi_j^2(x_j) = \varphi_j(x_j) = y_j,
    \]
    where we used $\varphi_j\varphi_i=0$ for $j\neq i$ and $\varphi_j^2=\varphi_j$. Hence $g(x)=(y_1,\ldots,y_n)$.

    \textit{Injective}: if $g(x)=0$, then $\varphi_i(x)=0$ for all $i$, so $x=\operatorname{id}_M(x)=\sum_i\varphi_i(x)=0$.

    Since $\operatorname{im}\varphi_i\simeq M_i$, we conclude $M\simeq\bigoplus_{i=1}^n M_i$.
\end{proof}

\begin{exercise}\leavevmode
    \begin{enumerate}
        \item Prove $\mathbb{Z}/m\mathbb{Z} \otimes_{\mathbb{Z}} \mathbb{Z}/n\mathbb{Z} \simeq \mathbb{Z}/(m,n)\mathbb{Z}$, where $m, n \in \mathbb{Z}$.
        \item Let $A$ be an abelian group, prove $A \otimes_{\mathbb{Z}} \mathbb{Z}/n\mathbb{Z} \simeq A/nA$.
        \item Let $A, B$ be two finitely generated abelian groups. What is $A \otimes_{\mathbb{Z}} B$?
    \end{enumerate}
\end{exercise}

\begin{proof}\leavevmode
    \begin{enumerate}
        \item We use the right-exact sequence $\mathbb{Z}\xrightarrow{\cdot\, n}\mathbb{Z}\to\mathbb{Z}/n\mathbb{Z}\to 0$. Tensoring with $\mathbb{Z}/m\mathbb{Z}$ over $\mathbb{Z}$ gives the right-exact sequence
        \[
            \mathbb{Z}/m\mathbb{Z} \xrightarrow{\cdot\, n} \mathbb{Z}/m\mathbb{Z} \to \mathbb{Z}/m\mathbb{Z}\otimes_\mathbb{Z}\mathbb{Z}/n\mathbb{Z}\to 0.
        \]
        Hence $\mathbb{Z}/m\mathbb{Z}\otimes_\mathbb{Z}\mathbb{Z}/n\mathbb{Z}\simeq (\mathbb{Z}/m\mathbb{Z})/n(\mathbb{Z}/m\mathbb{Z})$. Now $n(\mathbb{Z}/m\mathbb{Z})=(n\mathbb{Z}+m\mathbb{Z})/m\mathbb{Z}$. Since $n\mathbb{Z}+m\mathbb{Z}=(m,n)\mathbb{Z}$, we get
        \[
            \mathbb{Z}/m\mathbb{Z}\otimes_\mathbb{Z}\mathbb{Z}/n\mathbb{Z}\simeq \frac{\mathbb{Z}/m\mathbb{Z}}{(m,n)\mathbb{Z}/m\mathbb{Z}} \simeq \mathbb{Z}/(m,n)\mathbb{Z}.
        \]

        \item Similarly, tensor the right-exact sequence $\mathbb{Z}\xrightarrow{\cdot\, n}\mathbb{Z}\to\mathbb{Z}/n\mathbb{Z}\to 0$ with $A$:
        \[
            A \xrightarrow{\cdot\, n} A \to A\otimes_\mathbb{Z}\mathbb{Z}/n\mathbb{Z}\to 0.
        \]
        By right-exactness, $A\otimes_\mathbb{Z}\mathbb{Z}/n\mathbb{Z}\simeq A/nA$.

        \item By the fundamental theorem of finitely generated abelian groups,
        \[
            A\simeq \mathbb{Z}^r\oplus\bigoplus_{i=1}^s \mathbb{Z}/d_i\mathbb{Z}, \qquad B\simeq \mathbb{Z}^t\oplus\bigoplus_{j=1}^u \mathbb{Z}/e_j\mathbb{Z}.
        \]
        Since tensor product distributes over direct sums, using (1), (2), and $\mathbb{Z}\otimes_\mathbb{Z}M\simeq M$, we get
        \[
            A\otimes_\mathbb{Z}B \simeq \mathbb{Z}^{rt}\oplus \bigoplus_{j=1}^u(\mathbb{Z}/e_j\mathbb{Z})^r \oplus\bigoplus_{i=1}^s(\mathbb{Z}/d_i\mathbb{Z})^t \oplus \bigoplus_{i=1}^s\bigoplus_{j=1}^u\mathbb{Z}/(d_i,e_j)\mathbb{Z}.
        \]
    \end{enumerate}
\end{proof}

\begin{exercise}
    Let $\varphi: \mathbb{Z}/2\mathbb{Z} \to \mathbb{Z}/4\mathbb{Z}$ be the unique monomorphism. Prove that
    \[
    \operatorname{id} \otimes \varphi: \mathbb{Z}/2\mathbb{Z} \otimes_{\mathbb{Z}} \mathbb{Z}/2\mathbb{Z} \to \mathbb{Z}/2\mathbb{Z} \otimes_{\mathbb{Z}} \mathbb{Z}/4\mathbb{Z}
    \]
    is a zero homomorphism.
\end{exercise}

\begin{proof}
    The unique monomorphism $\varphi:\mathbb{Z}/2\mathbb{Z}\to\mathbb{Z}/4\mathbb{Z}$ is given by $\varphi(\bar 1)=\bar 2$. By the previous exercise, $\mathbb{Z}/2\mathbb{Z}\otimes_\mathbb{Z}\mathbb{Z}/2\mathbb{Z}\simeq\mathbb{Z}/2\mathbb{Z}$ and $\mathbb{Z}/2\mathbb{Z}\otimes_\mathbb{Z}\mathbb{Z}/4\mathbb{Z}\simeq\mathbb{Z}/2\mathbb{Z}$.

    It suffices to check on the generator $\bar 1\otimes\bar 1$ of $\mathbb{Z}/2\mathbb{Z}\otimes_\mathbb{Z}\mathbb{Z}/2\mathbb{Z}$:
    \[
        (\operatorname{id}\otimes\varphi)(\bar 1\otimes\bar 1) = \bar 1\otimes\varphi(\bar 1) = \bar 1\otimes\bar 2 = \bar 1\otimes 2\cdot\bar 1 = 2\cdot(\bar 1\otimes\bar 1) = (2\cdot\bar 1)\otimes\bar 1 = \bar 0\otimes\bar 1 = 0,
    \]
    where we used $2\cdot\bar 1 = \bar 0$ in $\mathbb{Z}/2\mathbb{Z}$. Since the generator maps to $0$, the map $\operatorname{id}\otimes\varphi$ is the zero homomorphism.

    This also shows that tensoring a monomorphism need not yield a monomorphism, i.e.\ the functor $\mathbb{Z}/2\mathbb{Z}\otimes_\mathbb{Z}-$ is not left-exact.
\end{proof}

\begin{exercise}
    Let $V_1$ and $V_2$ be two vector spaces over a field $F$ with dimension $n$ and $m$ respectively. Let $\{\epsilon_1, \ldots, \epsilon_n\}$ and $\{\eta_1, \ldots, \eta_m\}$ be bases of $V_1$ and $V_2$ respectively. Suppose that $\mathscr{A}$ (resp.\ $\mathscr{B}$) is an $F$-linear map of $V_1$ (resp.\ $V_2$), and it is represented by the matrix $A = (a_{ij})_{n \times n}$ (resp.\ $B = (b_{ij})_{m \times m}$) with respect to the above bases. Prove that
    \[
    \{\epsilon_i \otimes \eta_j \mid 1 \le i \le n,\, 1 \le j \le m\}
    \]
    is an $F$-basis of $V_1 \otimes_F V_2$. Determine the matrix representing the $F$-linear map $\mathscr{A} \otimes \mathscr{B}$ under this basis.
\end{exercise}

\begin{proof}
    \textbf{Step 1: $\{\epsilon_i\otimes\eta_j\}$ is a basis.}

    \textit{Spanning.} Every element of $V_1\otimes_F V_2$ is a finite sum of simple tensors $v\otimes w$. Writing $v=\sum_i a_i\epsilon_i$ and $w=\sum_j b_j\eta_j$, by bilinearity of the tensor product we get
    \[
        v\otimes w = \sum_{i,j} a_i b_j\,(\epsilon_i\otimes\eta_j).
    \]
    Hence $\{\epsilon_i\otimes\eta_j\}$ spans $V_1\otimes_F V_2$.

    \textit{Linear independence.} Suppose $\sum_{i,j}c_{ij}(\epsilon_i\otimes\eta_j)=0$. For each pair $(p,q)$, let $f_p\in V_1^*$ and $g_q\in V_2^*$ be the dual basis functionals: $f_p(\epsilon_i)=\delta_{pi}$ and $g_q(\eta_j)=\delta_{qj}$. The bilinear map $(v,w)\mapsto f_p(v)g_q(w)$ from $V_1\times V_2$ to $F$ induces (by the universal property of the tensor product) an $F$-linear map $\Phi_{pq}:V_1\otimes_F V_2\to F$ with $\Phi_{pq}(\epsilon_i\otimes\eta_j)=\delta_{pi}\delta_{qj}$. Applying $\Phi_{pq}$ to the relation $\sum_{i,j}c_{ij}(\epsilon_i\otimes\eta_j)=0$ yields $c_{pq}=0$. Since $(p,q)$ was arbitrary, all $c_{ij}=0$.

    Therefore $\{\epsilon_i\otimes\eta_j:1\le i\le n,\,1\le j\le m\}$ is an $F$-basis of $V_1\otimes_F V_2$, and $\dim_F(V_1\otimes_F V_2)=nm$.

    \medskip
    \textbf{Step 2: matrix of $\mathscr{A}\otimes\mathscr{B}$.}

    Order the basis lexicographically: $\epsilon_1\otimes\eta_1,\,\epsilon_1\otimes\eta_2,\,\ldots,\,\epsilon_1\otimes\eta_m,\,\epsilon_2\otimes\eta_1,\,\ldots,\,\epsilon_n\otimes\eta_m$. Since $\mathscr{A}(\epsilon_i)=\sum_{k=1}^n a_{ki}\epsilon_k$ and $\mathscr{B}(\eta_j)=\sum_{l=1}^m b_{lj}\eta_l$, we have
    \[
        (\mathscr{A}\otimes\mathscr{B})(\epsilon_i\otimes\eta_j) = \mathscr{A}(\epsilon_i)\otimes\mathscr{B}(\eta_j) = \sum_{k=1}^n\sum_{l=1}^m a_{ki}b_{lj}\,(\epsilon_k\otimes\eta_l).
    \]
    The entry in row $(k,l)$ and column $(i,j)$ of the representing matrix is $a_{ki}b_{lj}$. Written as an $nm\times nm$ block matrix, this is the \textbf{Kronecker product}
    \[
        A\otimes B = \begin{pmatrix} a_{11}B & a_{12}B & \cdots & a_{1n}B \\ a_{21}B & a_{22}B & \cdots & a_{2n}B \\ \vdots & \vdots & \ddots & \vdots \\ a_{n1}B & a_{n2}B & \cdots & a_{nn}B \end{pmatrix}.
    \]
\end{proof}

\end{document}
